% !TEX root = 第八章_数项级数(答案版).tex
\setcounter{chapter}{7}
\pagestyle{fancy}

\chapter{数项级数}


\section{级数收敛的定义及求和问题}

\begin{theorem}[\markstar\ 加括号]
    设 $\displaystyle\sum\limits_{n=0}^{\infty}a_n$ 是一收敛级数,如果把级数的项任意结合而不改变其先后次序,得新级数
    \[
        (a_1+\cdots+a_{k_1})+(a_{k_1+1}+\cdots+a_{k_2})+\cdots
        +(a_{k_{n-1}+1}+\cdots+a_{k_n})+\cdots,
    \]
    这里的正整数 $k_j$($j=1,2,\ldots$)满足 $k_1<k_2<\cdots$,那么新级数也收敛,且与原级数有相同的和.
\end{theorem}

\begin{theorem}[\markstar]
    若上式中同一括号中的项都有相同的符号,那么从加括号后的级数收敛便能推出原级数 $\displaystyle\sum\limits_{n=1}^{\infty}a_n$ 收敛,而且两者有相同的和.
\end{theorem}



\subsection{裂项相消法}

%例题8.1.1
\begin{example}
    设 $0<x<1$,求如下级数之和:
    \[
        \sum\limits_{n=0}^{\infty}\dfrac{x^{2n}}{1-x^{2n+1}}.
    \]
\end{example}
\exampleblank

\begin{solution}
    绝对收敛允许交换求和次序.将分母展开为几何级数,得
    \[
        \begin{aligned}
            I & =\sum\limits_{n=0}^{\infty}\sum\limits_{m=0}^{\infty}
            x^{2n+m(2n+1)}
            =\sum\limits_{m=0}^{\infty}x^m
            \sum\limits_{n=0}^{\infty}x^{2(m+1)n}                           \\
              & =\boxed{\sum\limits_{m=0}^{\infty}\dfrac{x^m}{1-x^{2m+2}}}.
        \end{aligned}
    \]
    按当前题面,该和通常不能再化为初等闭式;上式即其 Lambert 级数表示.
\end{solution}

%例题8.1.2
\begin{example}
    求下列级数之和:
    \begin{enumerate}
        \item $\displaystyle\sum\limits_{n=1}^{\infty}\dfrac{1}{(kn-i)(kn+j)}$,其中 $k=i+j$;
        \item $\displaystyle\sum\limits_{n=1}^{\infty}\dfrac{1}{n(n+1)(n+2)}$;
        \item $\displaystyle\sum\limits_{n=1}^{\infty}\dfrac{n-\sqrt{n^2-1}}{\sqrt{n(n+1)}}$;
        \item $\displaystyle\sum\limits_{n=1}^{\infty}(\sqrt{n}-2\sqrt{n+1}+\sqrt{n+2})$.
    \end{enumerate}
\end{example}
\exampleblank

\begin{solution}
    \begin{enumerate}
        \item 因 $k=i+j$,
              \[
                  \dfrac{1}{(kn-i)(kn+j)}
                  =\dfrac{1}{k}\left(\dfrac{1}{kn-i}-\dfrac{1}{k(n+1)-i}\right),
              \]
              故和为 $\boxed{\dfrac{1}{kj}}$.
        \item
              \[
                  \dfrac{1}{n(n+1)(n+2)}
                  =\dfrac{1}{2}\left[\dfrac{1}{n(n+1)}-\dfrac{1}{(n+1)(n+2)}\right],
              \]
              故和为 $\boxed{\dfrac{1}{4}}$.
        \item 通项等于
              \[
                  \sqrt{\dfrac{n}{n+1}}-\sqrt{\dfrac{n-1}{n}},
              \]
              所以部分和为 $\sqrt{\dfrac{N}{N+1}}$,级数和为 $\boxed1$.
        \item 前 $N$ 项之和为
              \[
                  1-\sqrt2-\sqrt{N+1}+\sqrt{N+2},
              \]
              故级数和为 $\boxed{1-\sqrt2}$.
    \end{enumerate}
\end{solution}

%例题8.1.3
\begin{example}
    求下列级数 $\displaystyle I=\sum\limits_{n=1}^{\infty}a_n$ 之和:
    \begin{enumerate}
        \item $\displaystyle a_n=\sin\dfrac{1}{2n}\cos\dfrac{3}{2n}$;
        \item $\displaystyle a_n=\arctan\dfrac{1}{2n^2}$;
        \item $\displaystyle a_n=\arctan\dfrac{2}{4n^2-4n+1}$;
        \item $\displaystyle a_n=3^{n-1}\sin^3\left(\dfrac{\theta}{3^n}\right)$.
    \end{enumerate}
\end{example}
\exampleblank

\begin{solution}
    \begin{enumerate}
        \item 由积化和差,
              \[
                  a_n=\dfrac{1}{2}\left(\sin\dfrac{2}{n}-\sin\dfrac{1}{n}\right)
                  \sim\dfrac{1}{2n}.
              \]
              因而按当前题面该级数发散,并不存在有限的``和'';这里的分母很可能漏写了指数,应核对原始资料.
        \item 利用反正切减法公式,
              \[
                  \arctan\dfrac{1}{2n^2}
                  =\arctan\dfrac{1}{2n-1}-\arctan\dfrac{1}{2n+1},
              \]
              故 $I=\boxed{\dfrac{\pi}{4}}$.
        \item 同理,
              \[
                  \arctan\dfrac{2}{4n^2-4n+1}
                  =\arctan\dfrac{1}{2n-2}-\arctan\dfrac{1}{2n},
              \]
              其中 $n=1$ 时把第一项理解为 $\dfrac{\pi}{2}$.故 $I=\boxed{\dfrac{\pi}{2}}$.
        \item 由 $\sin3u=3\sin u-4\sin^3u$,
              \[
                  3^{n-1}\sin^3\dfrac{\theta}{3^n}
                  =\dfrac{1}{4}\left(3^n\sin\dfrac{\theta}{3^n}
                  -3^{n-1}\sin\dfrac{\theta}{3^{n-1}}\right).
              \]
              裂项相消并令 $N\to\infty$,得到
              \[
                  I=\boxed{\dfrac{\theta-\sin\theta}{4}}.
              \]
    \end{enumerate}
\end{solution}

%例题8.1.4
\begin{example}
    已知级数 $\displaystyle\sum\limits_{n=1}^{\infty}a_n=A$,$\displaystyle\lim\limits_{n\to\infty}na_n=0$,求级数
    \[
        \sum\limits_{n=1}^{\infty}n(a_n-a_{n+1})
    \]
    之和.
\end{example}
\exampleblank

\begin{solution}
    前 $N$ 项满足
    \[
        \begin{aligned}
            \sum\limits_{n=1}^{N}n(a_n-a_{n+1})
             & =\sum\limits_{n=1}^{N}na_n-
            \sum\limits_{n=2}^{N+1}(n-1)a_n        \\
             & =\sum\limits_{n=1}^{N}a_n-Na_{N+1}.
        \end{aligned}
    \]
    由 $na_n\to0$,右端趋于 $A$,故所求和为 $\boxed A$.
\end{solution}

%例题8.1.5
\begin{example}
    设 $a_n>0$ 且 $\displaystyle\sum\limits_{n=1}^{\infty}a_n=S$,若记
    \[
        R_n=\sum\limits_{k=n+1}^{\infty}a_k,
    \]
    则
    \[
        I=\sum\limits_{n=1}^{\infty}\dfrac{a_n}{\sqrt{R_{n-1}}+\sqrt{R_n}}=\sqrt{S}.
    \]
\end{example}
\exampleblank

\begin{solution}
    由 $a_n=R_{n-1}-R_n$,
    \[
        \dfrac{a_n}{\sqrt{R_{n-1}}+\sqrt{R_n}}
        =\sqrt{R_{n-1}}-\sqrt{R_n}.
    \]
    前 $N$ 项之和为 $\sqrt{R_0}-\sqrt{R_N}$.又 $R_0=S$ 且 $R_N\to0$,故
    \[
        I=\boxed{\sqrt S}.
    \]
\end{solution}

%例题8.1.6
\begin{example}
    设数列 $\{a_n\}$ 满足
    \[
        \sum\limits_{n=1}^{\infty}(a_1+1)(a_2+1)\cdots(a_n+1)=l,
    \]
    其中 $l>0$ 或 $l=+\infty$.试证明:
    \[
        \sum\limits_{n=1}^{\infty}
        \dfrac{a_n}{(a_1+1)(a_2+1)\cdots(a_n+1)}
        =1-\dfrac{1}{l}.
    \]
\end{example}
\exampleblank

\begin{solution}
    当前题面不成立.令
    \[
        P_n=(1+a_1)(1+a_2)\cdots(1+a_n),\qquad P_0=1,
    \]
    则恒有
    \[
        \dfrac{a_n}{P_n}=\dfrac{1}{P_{n-1}}-\dfrac{1}{P_n}.
    \]
    因此若正确条件是 $P_n\to l$,结论才由裂项相消得到
    \[
        \sum\limits_{n=1}^{\infty}\dfrac{a_n}{P_n}=1-\dfrac{1}{l}.
    \]
    而原题写的是 $\displaystyle \sum P_n=l$.例如取 $a_n=0$,则 $l=+\infty$,但左端为 $0$,右端为 $1$.故条件中的求和号应核对原始资料.
\end{solution}

%例题8.1.7
\begin{example}
    设 $a_1>1$,$a_{n+1}=a_n^2-a_n+1$,则
    \[
        I=\sum\limits_{n=1}^{\infty}\dfrac{1}{a_n}=\dfrac{1}{a_1-1}.
    \]
\end{example}
\exampleblank

\begin{solution}
    递推式给出
    \[
        a_{n+1}-1=a_n(a_n-1),
    \]
    因而
    \[
        \dfrac{1}{a_n}
        =\dfrac{1}{a_n-1}-\dfrac{1}{a_{n+1}-1}.
    \]
    又 $a_n>1$ 且严格递增,故 $a_n\to+\infty$.裂项相消即得
    \[
        I=\boxed{\dfrac{1}{a_1-1}}.
    \]
\end{solution}



\subsection{子列方法}

易知,若 $\{S_{2n}\}$ 与 $\{S_{2n+1}\}$ 有相同极限 $S$,则 $\displaystyle\lim\limits_{n\to\infty}S_n=S$.因此若通项 $a_n\to0$,则可由 $S_{2n}\to S$ 得到 $S_{2n+1}\to S$,进而 $\displaystyle\sum\limits_{n=1}^{\infty}a_n=S$.更一般地,若 $a_n\to0$ 且 $\{S_{pn}\}\to S$($p$ 是某个正整数),则 $\displaystyle\sum\limits_{n=1}^{\infty}a_n=S$.我们称这种方法为子列方法.

%例题8.1.8
\begin{example}
    计算
    \[
        1+\dfrac{1}{2}+\left(\dfrac{1}{3}-1\right)+\dfrac{1}{4}+\dfrac{1}{5}
        +\left(\dfrac{1}{6}-\dfrac{1}{2}\right)+\dfrac{1}{7}+\dfrac{1}{8}
        +\left(\dfrac{1}{9}-\dfrac{1}{3}\right)+\cdots.
    \]
\end{example}
\exampleblank

\begin{solution}
    取到第 $3N$ 项,括号中的减项恰为 $1,\dfrac{1}{2},\ldots,\dfrac{1}{N}$,故部分和为
    \[
        H_{3N}-H_N,
        \qquad H_N=\sum\limits_{k=1}^{N}\dfrac{1}{k}.
    \]
    由 $H_N=\ln N+\gamma+o(1)$,所求和为 $\boxed{\ln3}$.
\end{solution}

%例题8.1.9
\begin{example}
    计算
    \[
        1-\dfrac{1}{2}-\dfrac{1}{4}+\dfrac{1}{3}-\dfrac{1}{6}-\dfrac{1}{8}
        +\cdots+\dfrac{1}{2n-1}-\dfrac{1}{4n-2}-\dfrac{1}{4n}+\cdots.
    \]
\end{example}
\exampleblank

\begin{solution}
    每组三项之和为
    \[
        \dfrac{1}{2n-1}-\dfrac{1}{4n-2}-\dfrac{1}{4n}
        =\dfrac{1}{2}\left(\dfrac{1}{2n-1}-\dfrac{1}{2n}\right).
    \]
    因此所求和为交错调和级数之和的一半,即 $\boxed{\dfrac{1}{2}\ln2}$.
\end{solution}

%例题8.1.10
\begin{example}
    将级数 $\displaystyle1-\dfrac{1}{2}+\dfrac{1}{3}-\dfrac{1}{4}+\dfrac{1}{5}-\cdots$ 的各项重新安排,使先依次出现 $p$ 个正项,再出现 $q$ 个负项,然后如此交替.试证新级数的和为
    \[
        \ln2+\dfrac{1}{2}\ln\dfrac{p}{q}.
    \]
\end{example}
\exampleblank

\begin{solution}
    记重排后前 $N$ 个完整组的和为 $T_N$,则
    \[
        T_N=\sum\limits_{k=1}^{pN}\dfrac{1}{2k-1}
        -\sum\limits_{k=1}^{qN}\dfrac{1}{2k}.
    \]
    利用
    \[
        \sum\limits_{k=1}^{m}\dfrac{1}{2k-1}=H_{2m}-\dfrac{1}{2}H_m,
        \qquad
        \sum\limits_{k=1}^{m}\dfrac{1}{2k}=\dfrac{1}{2}H_m,
    \]
    可得
    \[
        T_N=H_{2pN}-\dfrac{1}{2}H_{pN}-\dfrac{1}{2}H_{qN}
        \longrightarrow\ln2+\dfrac{1}{2}\ln\dfrac{p}{q}.
    \]
    一个组内尚未取完的项总和趋于零,故所有部分和具有同一极限.
\end{solution}

%例题8.1.11
\begin{example}
    证明:若删去调和级数中所有分母含有数字 $9$ 的项,则新级数收敛,且和小于 $80$.
\end{example}
\exampleblank

\begin{solution}
    按十进制位数分组.恰有 $k$ 位且分母中不含数字 $9$ 的正整数至多有
    \[
        8\cdot9^{k-1}
    \]
    个,每个均不小于 $10^{k-1}$.因此第 $k$ 组倒数之和严格小于
    \[
        8\left(\dfrac{9}{10}\right)^{k-1}.
    \]
    故新级数收敛,并且其和小于
    \[
        8\sum\limits_{k=1}^{\infty}\left(\dfrac{9}{10}\right)^{k-1}
        =\boxed{80}.
    \]
\end{solution}



\subsection{发散级数}

%例题8.1.12
\begin{example}
    试证明下列级数的发散性:
    \begin{enumerate}
        \item $\displaystyle\sum\limits_{n=1}^{\infty}(n^2+2)\ln\left(\dfrac{n^2+1}{n^2}\right)$;
        \item $\displaystyle\sum\limits_{n=1}^{\infty}\left(\dfrac{2n^2-3}{2n^2+1}\right)^{n^2}$;
        \item $\displaystyle\sum\limits_{n=1}^{\infty}\cos(n\alpha)$;
        \item $\displaystyle\sum\limits_{n=1}^{\infty}\dfrac{1}{n}$.
    \end{enumerate}
\end{example}
\exampleblank

\begin{solution}
    \begin{enumerate}
        \item 因
              \[
                  (n^2+2)\ln\left(1+\dfrac{1}{n^2}\right)\longrightarrow1,
              \]
              通项不趋于零.
        \item
              \[
                  \left(1-\dfrac{4}{2n^2+1}\right)^{n^2}\longrightarrow e^{-2}\neq0.
              \]
        \item 若 $\cos(n\alpha)\to0$,则其子列 $\cos(2n\alpha)\to0$;但恒等式
              $\cos(2n\alpha)=2\cos^2(n\alpha)-1$ 又给出该子列趋于 $-1$,矛盾.
        \item 由分组估计
              \[
                  \sum\limits_{n=2^{m-1}+1}^{2^m}\dfrac{1}{n}\geq\dfrac{1}{2},
              \]
              调和级数发散.
    \end{enumerate}
\end{solution}

%例题8.1.13
\begin{example}
    设 $a_n>0$,$\displaystyle\sum\limits_{n=1}^{\infty}a_n=+\infty$,试证:
    \[
        \sum\limits_{n=1}^{\infty}\dfrac{a_n}{S_n}
    \]
    发散.
\end{example}
\exampleblank

\begin{solution}
    记 $\displaystyle S_n=\sum\limits_{k=1}^{n}a_k$.若 $a_n/S_n$ 不趋于零,则结论显然.否则最终有 $a_n/S_n\leq\dfrac{1}{2}$,即 $a_n/S_{n-1}\leq1$.此时
    \[
        \ln\dfrac{S_n}{S_{n-1}}
        =\ln\left(1+\dfrac{a_n}{S_{n-1}}\right)
        \leq\dfrac{2a_n}{S_n}.
    \]
    求和后得到
    \[
        \sum\limits_{n=N}^{m}\dfrac{a_n}{S_n}
        \geq\dfrac{1}{2}\ln\dfrac{S_m}{S_{N-1}}\longrightarrow+\infty,
    \]
    故原级数发散.
\end{solution}

%例题8.1.14
\begin{example}
    如果
    \[
        \lim\limits_{n\to\infty}a_{n+1}=0,
        \quad
        \lim\limits_{n\to\infty}(a_{n+1}+a_{n+2})=0,
        \quad\ldots,\quad
        \lim\limits_{n\to\infty}(a_{n+1}+\cdots+a_{n+p})=0,
    \]
    试问 $\displaystyle\sum\limits_{n=1}^{\infty}a_n$ 是否一定收敛.
\end{example}
\exampleblank

\begin{solution}
    不一定.事实上题设第一式已经是 $a_{n+1}\to0$,其余每个固定长度的和也随之趋于零,并未增加足以保证级数收敛的条件.取
    \[
        a_n=\dfrac{1}{n}
    \]
    即可满足全部极限条件,但调和级数发散.
\end{solution}

%例题8.1.15
\begin{example}
    设 $\{a_n\}$ 是递减正数列,且 $\displaystyle\sum\limits_{n=1}^{\infty}a_n=+\infty$,证明:
    \[
        \lim\limits_{n\to\infty}
        \dfrac{a_1+a_3+\cdots+a_{2n-1}}{a_2+a_4+\cdots+a_{2n}}=1.
    \]
\end{example}
\exampleblank

\begin{solution}
    记
    \[
        O_n=\sum\limits_{k=1}^{n}a_{2k-1},\qquad
        E_n=\sum\limits_{k=1}^{n}a_{2k}.
    \]
    单调性给出 $O_n\geq E_n$,并且
    \[
        0\leq O_n-E_n
        =\sum\limits_{k=1}^{n}(a_{2k-1}-a_{2k})
        \leq\sum\limits_{k=1}^{n}(a_{2k-1}-a_{2k+1})
        \leq a_1.
    \]
    又 $O_n+E_n\to+\infty$,故 $E_n\to+\infty$,从而
    \[
        \dfrac{O_n}{E_n}=1+\dfrac{O_n-E_n}{E_n}\longrightarrow\boxed1.
    \]
\end{solution}

%例题8.1.16
\begin{example}
    证明下列级数发散:
    \begin{enumerate}
        \item $\displaystyle1-\dfrac{1}{\sqrt2}+\dfrac{1}{3}-\dfrac{1}{\sqrt4}+\dfrac{1}{5}-\dfrac{1}{\sqrt6}+\cdots$;
        \item $\displaystyle1+\dfrac{1}{\sqrt3}-\dfrac{1}{\sqrt2}+\dfrac{1}{\sqrt5}+\dfrac{1}{\sqrt7}-\dfrac{1}{\sqrt4}+\cdots$.
    \end{enumerate}
\end{example}
\exampleblank

\begin{solution}
    \begin{enumerate}
        \item 两两分组后第 $k$ 组为
              \[
                  \dfrac{1}{2k-1}-\dfrac{1}{\sqrt{2k}}.
              \]
              当 $k$ 充分大时它不超过 $-c/\sqrt k$,其中 $c>0$,故分组部分和趋于 $-\infty$.
        \item 每组三项为
              \[
                  \dfrac{1}{\sqrt{4k-3}}+\dfrac{1}{\sqrt{4k-1}}-\dfrac{1}{\sqrt{2k}}.
              \]
              乘以 $\sqrt k$ 后趋于 $1-\dfrac{1}{\sqrt2}>0$,故与 $\displaystyle \sum\limits k^{-1/2}$ 比较可知分组部分和趋于 $+\infty$.
    \end{enumerate}
\end{solution}


\newpage

\section{正项级数}

判断级数 $\displaystyle\sum a_n$ 的敛散性,通常有如下方法:
\begin{enumerate}
    \item 若通项 $a_n\not\to0$,则级数发散;
    \item 考虑部分和 $S_n$ 是否有界,有界则收敛,无界则发散;
    \item Cauchy 收敛准则;
    \item 比较判别法.根据被比较的级数不同,又可细分为 Cauchy 凝聚判别法,Cauchy 积分判别法,Cauchy 根式判别法,d'Alembert 比值判别法,Rabbe 判别法,Gauss 判别法等.
\end{enumerate}



\subsection{一般比较判别}


\methodsection{(A) 通项大小比较法}

\begin{theorem}[\marktwostar]
    设有两个正项级数
    \[
        (A)\sim\sum\limits_{n=1}^{\infty}a_n,
        \qquad
        (B)\sim\sum\limits_{n=1}^{\infty}b_n,
    \]
    满足 $0<a_n\leq b_n$($n\geq n_0$).则:
    \begin{enumerate}
        \item 若 $(B)$ 收敛,则 $(A)$ 收敛;
        \item 若 $(A)$ 发散,则 $(B)$ 发散.
    \end{enumerate}
\end{theorem}

\begin{remark}
    若有 $0<d_n<M$,使得 $0<a_n\leq d_nb_n$,则上述第一条仍成立.
\end{remark}

\begin{corollary}[\marktwostar]
    设 $a_n,b_n>0$,且
    \[
        \dfrac{a_{n+1}}{a_n}\leq\dfrac{b_{n+1}}{b_n},
    \]
    则由 $\displaystyle\sum\limits_{n=1}^{\infty}b_n$ 收敛可导出 $\displaystyle\sum\limits_{n=1}^{\infty}a_n$ 收敛.
\end{corollary}

\begin{remark}
    因此定理 8.2.1及推论 8.2.1的核心在于利用不等式估计与已知敛散性的级数进而比较,这就要求不等式基础了.
\end{remark}

%例题8.2.1
\begin{example}
    判断级数 $\displaystyle I=\sum\limits_{n=3}^{\infty}a_n$ 的敛散性:
    \begin{enumerate}
        \item $\displaystyle a_n=\dfrac{1}{\sqrt[r]{\ln n}}\quad(r>0)$;
        \item $\displaystyle a_n=\dfrac{1}{(\ln n)^{\ln n}}$;
        \item $\displaystyle a_n=\dfrac{1}{(\ln\ln n)^{\ln n}}$;
        \item $\displaystyle a_n=\dfrac{1}{n^{1+\frac{1}{\ln\ln n}}}$.
    \end{enumerate}
\end{example}
\exampleblank

\begin{solution}
    \begin{enumerate}
        \item 对充分大的 $n$,$(\ln n)^{1/r}<n$,故 $a_n>1/n$,级数发散.
        \item
              \[
                  a_n=e^{-(\ln n)(\ln\ln n)}=n^{-\ln\ln n}.
              \]
              最终有 $a_n\leq n^{-2}$,故收敛.
        \item 同理 $a_n=n^{-\ln\ln\ln n}$,最终也有 $a_n\leq n^{-2}$,故收敛.
        \item 由 Cauchy 凝聚判别,凝聚后的通项为
              \[
                  2^na_{2^n}=\exp\left(-\dfrac{n\ln2}{\ln(n\ln2)}\right).
              \]
              该量最终小于 $n^{-2}$,故原级数收敛.
    \end{enumerate}
\end{solution}

%例题8.2.2
\begin{example}
    判断级数 $\displaystyle I=\sum\limits_{n=1}^{\infty}a_n$ 的敛散性:
    \begin{enumerate}
        \item $\displaystyle a_n=\dfrac{1}{\ln(n!)}\quad(n\geq2)$;
        \item $\displaystyle a_n=\dfrac{\ln(n!)}{n^p}$;
        \item $\displaystyle a_n=\dfrac{\sum\limits_{k=1}^n\ln^2k}{n^p}$;
        \item $\displaystyle a_n=\left[\dfrac{(2n-1)!!}{(2n)!!}\right]^2$.
    \end{enumerate}
\end{example}
\exampleblank

\begin{solution}
    由 Stirling 公式与积分比较,
    \[
        \ln(n!)\sim n\ln n,
        \qquad
        \sum\limits_{k=1}^{n}\ln^2k\sim n\ln^2n.
    \]
    因此 (1) 与 $\displaystyle \sum 1/(n\ln n)$ 同敛散,故发散;(2) 与
    $\displaystyle \sum (\ln n)/n^{p-1}$ 同敛散,恰在 $p>2$ 时收敛;(3) 与
    $\displaystyle \sum (\ln^2n)/n^{p-1}$ 同敛散,也恰在 $p>2$ 时收敛.

    对 (4),Wallis 公式给出
    \[
        \dfrac{(2n-1)!!}{(2n)!!}\sim\dfrac{1}{\sqrt{\pi n}},
    \]
    故通项与 $1/(\pi n)$ 等价,级数发散.
\end{solution}

%例题8.2.3
\begin{example}
    试证明下列命题:
    \begin{enumerate}
        \item 设 $a_n>0$,若 $\displaystyle\sum\limits_{n=1}^{\infty}a_n$ 收敛,则
              \[
                  \sum\limits_{n=1}^{\infty}\dfrac{\sqrt{a_n}}{n^p}
              \]
              在 $p>\dfrac{1}{2}$ 时收敛;
        \item 若 $\displaystyle I=\sum\limits_{n=1}^{\infty}a_n$ 是正项收敛级数,则 $\displaystyle J=\sum\limits_{n=1}^{\infty}\sqrt{a_na_{n+1}}$ 收敛,反之不然.
    \end{enumerate}
\end{example}
\exampleblank

\begin{solution}
    \begin{enumerate}
        \item 由 Cauchy--Schwarz 不等式,
              \[
                  \sum\limits_{n=1}^{\infty}\dfrac{\sqrt{a_n}}{n^p}
                  \leq
                  \left(\sum\limits_{n=1}^{\infty}a_n\right)^{1/2}
                  \left(\sum\limits_{n=1}^{\infty}\dfrac{1}{n^{2p}}\right)^{1/2}<\infty
              \]
              当且仅当这里所用的 $p$ 满足 $p>1/2$.
        \item 由 $2\sqrt{a_na_{n+1}}\leq a_n+a_{n+1}$,$I$ 收敛蕴含 $J$ 收敛.反之不成立:取
              \[
                  a_{2k-1}=1,qquad a_{2k}=\dfrac{1}{k^4}.
              \]
              则 $\displaystyle \sum a_n$ 发散,而 $\displaystyle \sum\sqrt{a_na_{n+1}}$ 可由 $\displaystyle \sum k^{-2}$ 控制,因而收敛.
    \end{enumerate}
\end{solution}

%例题8.2.4
\begin{example}
    设 $\displaystyle\sum\limits_{n=1}^{\infty}a_n$ 是正项发散级数,试判断下列级数的敛散性:
    \begin{enumerate}
        \item $\displaystyle\sum\limits_{n=1}^{\infty}\dfrac{a_n}{1+a_n}$;
        \item $\displaystyle\sum\limits_{n=1}^{\infty}\dfrac{a_n}{1+na_n}$;
        \item $\displaystyle\sum\limits_{n=1}^{\infty}\dfrac{a_n}{1+n^2a_n}$;
        \item $\displaystyle\sum\limits_{n=1}^{\infty}\dfrac{a_n}{1+a_n^2}$.
    \end{enumerate}
\end{example}
\exampleblank

\begin{solution}
    \begin{enumerate}
        \item 恒发散.若有无穷多个 $a_n\geq1$,相应通项至少为 $1/2$;否则最终有 $a_n<1$,相应通项至少为 $a_n/2$.
        \item 不能确定.取 $a_n=1/n$ 时所得级数发散;若 $a_{2^k}=1$,其余项取 $a_n=1/n^2$,则 $\displaystyle \sum a_n$ 发散,而所得级数收敛.
        \item 恒收敛,因为
              \[
                  0<\dfrac{a_n}{1+n^2a_n}\leq\dfrac{1}{n^2}.
              \]
        \item 不能确定.取 $a_n=1/n$ 时发散;取 $a_n=2^n$ 时通项小于 $2^{-n}$,故收敛.
    \end{enumerate}
\end{solution}

%例题8.2.5
\begin{example}
    设 $a_n\neq0$,且 $a_n\to a\neq0$,则
    \[
        \sum\limits_{n=1}^{\infty}|a_{n+1}-a_n|
    \]
    与
    \[
        \sum\limits_{n=1}^{\infty}\left|\dfrac{1}{a_{n+1}}-\dfrac{1}{a_m}\right|
    \]
    同敛散.
\end{example}
\exampleblank

\begin{solution}
    题面第二个级数中的 $a_m$ 应为 $a_n$.由于 $a_n\to a\neq0$,存在 $N$ 与常数 $0<c<C$,使 $n\geq N$ 时
    \[
        c\leq|a_n|\leq C.
    \]
    于是
    \[
        \dfrac{1}{C^2}|a_{n+1}-a_n|
        \leq
        \left|\dfrac{1}{a_{n+1}}-\dfrac{1}{a_n}\right|
        \leq\dfrac{1}{c^2}|a_{n+1}-a_n|.
    \]
    有限个首项不影响敛散性,故两级数同敛散.原题下标需要核对.
\end{solution}

%例题8.2.6
\begin{example}[\markstar\ Sapagof 判别法]
    设正数数列 $\{a_n\}$ 单调递减,则 $\displaystyle\lim\limits_{n\to\infty}a_n=0$ 的充分必要条件是正项级数
    \[
        \sum\limits_{n=1}^{\infty}\left(1-\dfrac{a_{n+1}}{a_n}\right)
    \]
    发散.
\end{example}
\exampleblank

\begin{solution}
    令 $r_n=a_{n+1}/a_n\in(0,1]$,则
    \[
        \dfrac{a_{N+1}}{a_1}=\prod\limits_{n=1}^{N}r_n.
    \]
    正项无穷乘积的基本判别说明
    \[
        \prod\limits_{n=1}^{\infty}r_n>0
        \quad\Longleftrightarrow\quad
        \sum\limits_{n=1}^{\infty}(1-r_n)<\infty;
    \]
    这也可由 $1-r\leq-\ln r$ 以及 $-\ln r\leq2(1-r)$($r$ 充分接近 $1$)直接得到.因此 $a_n$ 的极限为正数当且仅当题中级数收敛,取否命题即得结论.
\end{solution}

\begin{remark}
    由此可得:
    \begin{enumerate}
        \item 设 $\{a_n\}$ 为单调递增的正数数列,则该数列与级数 $\displaystyle\sum\limits_{n=1}^{\infty}\left(1-\dfrac{a_n}{a_{n+1}}\right)$ 同敛散;
        \item 若正项级数 $\displaystyle\sum\limits_{n=1}^{\infty}a_n$ 的部分和数列为 $\{S_n\}$,则 $\displaystyle\sum\limits_{n=1}^{\infty}a_n$ 与 $\displaystyle\sum\limits_{n=1}^{\infty}\dfrac{a_n}{S_n}$ 同敛散.
    \end{enumerate}
\end{remark}

\begin{proposition}[\markstar]
    \begin{enumerate}
        \item 设正项级数 $\displaystyle\sum\limits_{n=1}^{\infty}a_n$ 收敛,记 $\displaystyle R_n=\sum\limits_{k=n}^{\infty}a_k$,则级数 $\displaystyle\sum\limits_{n=1}^{\infty}\dfrac{a_n}{R_n^p}$ 在 $p<1$ 时收敛,$p\geq1$ 时发散;
        \item 设正项级数 $\displaystyle\sum\limits_{n=1}^{\infty}a_n$ 发散,记 $\displaystyle S_n=\sum\limits_{k=1}^na_k$,则级数 $\displaystyle\sum\limits_{n=1}^{\infty}\dfrac{a_n}{S_n^p}$ 在 $p\leq1$ 时发散,$p>1$ 时收敛.
    \end{enumerate}
\end{proposition}

\begin{remark}
    由证明过程可知,任意正项级数 $\displaystyle\sum\limits_{n=1}^{\infty}a_n$,当 $p>1$ 时,必有 $\displaystyle\sum\limits_{n=1}^{\infty}\dfrac{a_n}{S_n^p}$ 收敛.并且由上述命题,我们可知没有最小的发散级数,也没有最大的发散级数,因此没有所谓的``最完美''的判别法.
\end{remark}

%例题8.2.7
\begin{example}[\markstar]
    证明如下问题:
    \begin{enumerate}
        \item 设正项级数 $\displaystyle\sum\limits_{n=1}^{\infty}a_n$ 发散,则 $\displaystyle\sum\limits_{n=1}^{\infty}\dfrac{a_n}{S_n}$ 发散,$\displaystyle\sum\limits_{n=1}^{\infty}\dfrac{a_n}{S_n^2}$ 收敛;
        \item 设正项级数 $\displaystyle\sum\limits_{n=1}^{\infty}a_n$ 收敛,则 $\displaystyle\sum\limits_{n=1}^{\infty}\dfrac{a_n}{R_n}$ 发散,$\displaystyle\sum\limits_{n=1}^{\infty}\dfrac{a_n}{\sqrt{R_n}}$ 收敛;
        \item 设正项级数 $\displaystyle\sum\limits_{n=1}^{\infty}a_n$ 发散,则存在发散的正项级数 $\displaystyle\sum\limits_{n=1}^{\infty}b_n$,使得 $\displaystyle\sum\limits_{n=1}^{\infty}\dfrac{a_n}{b_n}=0$;
        \item 设正项级数 $\displaystyle\sum\limits_{n=1}^{\infty}a_n$ 收敛,则存在收敛的正项级数 $\displaystyle\sum\limits_{n=1}^{\infty}b_n$,使得 $\displaystyle\sum\limits_{n=1}^{\infty}\dfrac{a_n}{b_n}=0$.
    \end{enumerate}
\end{example}
\exampleblank

\begin{solution}
    \begin{enumerate}
        \item 记 $\displaystyle S_n=\sum_{k=1}^{n}a_k$.与前题相同,$\displaystyle \sum a_n/S_n$ 发散;又
              \[
                  \dfrac{a_n}{S_n^2}
                  \leq\dfrac{a_n}{S_{n-1}S_n}
                  =\dfrac{1}{S_{n-1}}-\dfrac{1}{S_n},
              \]
              故第二个级数收敛.
        \item 按本章前文的定义 $\displaystyle R_n=\sum_{k=n+1}^{\infty}a_k$,当前题面不成立.例如取 $a_n=2^{-2^n}$,则 $a_n/\sqrt{R_n}$ 的通项甚至不趋于零.若分母应为 $R_{n-1}$,则
              \[
                  \dfrac{a_n}{R_{n-1}}=1-\dfrac{R_n}{R_{n-1}},
                  \qquad
                  \dfrac{a_n}{\sqrt{R_{n-1}}}
                  \leq2(\sqrt{R_{n-1}}-\sqrt{R_n}),
              \]
              分别给出发散与收敛.故此处下标需要核对原始资料.
        \item 当前等式不可能成立,因为各项均为正数.若原结论是
              $a_n/b_n\to0$,可取 $b_n=na_n$;则 $\displaystyle \sum b_n$ 发散且 $a_n/b_n=1/n\to0$.
        \item 此处同样应为极限 $a_n/b_n\to0$.可取整数
              $1=N_1<N_2<\cdots$ 使
              \[
                  \sum\limits_{n=N_j}^{\infty}a_n<\dfrac{2^{-j}}{j},
              \]
              并在 $N_j\leq n<N_{j+1}$ 上置 $b_n=ja_n$.于是 $b_n>0$,$a_n/b_n=1/j\to0$,且
              \[
                  \sum\limits_{n=1}^{\infty}b_n
                  \leq\sum\limits_{j=1}^{\infty}j\sum\limits_{n=N_j}^{\infty}a_n
                  <\sum\limits_{j=1}^{\infty}2^{-j}<\infty.
              \]
    \end{enumerate}
\end{solution}

%例题8.2.8
\begin{example}[\markstar]
    举例说明下列命题中的存在性:
    \begin{enumerate}
        \item 设 $\{b_n\}$ 是趋于 $+\infty$ 的正数列,则必存在收敛的正项级数 $\displaystyle\sum\limits_{n=1}^{\infty}a_n$,使得 $\displaystyle\sum\limits_{n=1}^{\infty}a_nb_n=+\infty$;
        \item 设 $\displaystyle\sum\limits_{n=1}^{\infty}a_n$ 是收敛的正项级数,则存在正数列 $\{\lambda_n\}$ 满足 $\displaystyle\lim\limits_{n\to\infty}\lambda_n=+\infty$ 且 $\displaystyle\sum\limits_{n=1}^{\infty}\lambda_na_n<+\infty$;
        \item 存在 $\{a_n\},\{b_n\}$ 是递减趋于零的数列,且 $\displaystyle\sum\limits_{n=1}^{\infty}a_n,\sum\limits_{n=1}^{\infty}b_n$ 都发散,但 $\displaystyle\sum\limits_{n=1}^{\infty}\min\{a_n,b_n\}$ 收敛;
        \item 设 $\{a_n\}$ 是趋于零的正数列,且 $\displaystyle\sum\limits_{n=1}^{\infty}a_n=+\infty$,则对 $C>0$,存在子列 $\{n_k\}$,使得 $\displaystyle\sum\limits_{k=1}^{\infty}a_{n_k}=C$.
    \end{enumerate}
\end{example}
\exampleblank

\begin{solution}
    \begin{enumerate}
        \item 取子列 $n_k$ 使 $b_{n_k}\geq2^k$,并令
              \[
                  a_{n_k}=2^{-k},qquad
                  a_n=\dfrac{2^{-n}}{1+b_n}\quad(n\neq n_k).
              \]
              则 $\displaystyle \sum a_n<\infty$,而 $\displaystyle \sum a_nb_n\geq\sum 1=+\infty$.
        \item 取 $N_j$ 使尾和 $\displaystyle \sum_{n=N_j}^{\infty}a_n<2^{-j}/j$,并在
              $N_j\leq n<N_{j+1}$ 上令 $\lambda_n=j$.则 $\lambda_n\to\infty$,且
              $\displaystyle \sum\lambda_na_n\leq\sum2^{-j}<\infty$.
        \item 可递归构造相邻区间 $I_k$ 与递减高度 $\alpha_k,\beta_k$:在奇数块令
              $a_n=\alpha_k,b_n=\beta_k$ 且 $|I_k|\alpha_k\geq1$,$|I_k|\beta_k\leq2^{-k}$;在偶数块交换两者的角色.每次先把新高度取得比上一块两高度都小,再把块长取得足够大,即可同时满足这些条件.所得两列均递减趋零,$\displaystyle \sum a_n$ 与 $\displaystyle \sum b_n$ 因交替出现质量至少为 $1$ 的块而发散,而
              \[
                  \sum\limits_n\min\{a_n,b_n\}\leq\sum\limits_k2^{-k}<\infty.
              \]
        \item 采用贪心选取:依次考察 $a_n$,只要加入后不超过 $C$ 就选取该项.所得子级数和不超过 $C$.若其和为 $L<C$,由 $a_n\to0$,充分靠后的每项均小于 $C-L$;又原级数发散,之后必还能选入一项,矛盾.因此所选子列之和恰为 $C$.
    \end{enumerate}
\end{solution}

%例题8.2.9
\begin{example}
    设 $0<p_1<p_2<\cdots<p_n<\cdots$,求证:
    \[
        \sum\limits_{n=1}^{\infty}\dfrac{1}{p_n}\text{ 收敛}
        \quad\Longleftrightarrow\quad
        \sum\limits_{n=1}^{\infty}\dfrac{n}{p_1+p_2+\cdots+p_n}\text{ 收敛}.
    \]
\end{example}
\exampleblank

\begin{solution}
    记 $P_n=p_1+cdots+p_n$.由 $P_n\leq np_n$,
    \[
        \dfrac{1}{p_n}\leq\dfrac{n}{P_n},
    \]
    故右侧收敛蕴含左侧收敛.反之,由单调性,
    \[
        P_n\geq\sum\limits_{k=\lfloor n/2\rfloor}^{n}p_k
        \geq\dfrac{n}{2}p_{\lfloor n/2\rfloor},
    \]
    从而
    \[
        \dfrac{n}{P_n}\leq\dfrac{2}{p_{\lfloor n/2\rfloor}}.
    \]
    右端每个下标至多重复三次,故可与 $\displaystyle \sum1/p_n$ 比较,得到反向蕴含.
\end{solution}

\begin{remark}
    事实上,必要性证明可舍去单调性条件.
\end{remark}

%例题8.2.10
\begin{example}
    证明下列不等式:
    \begin{enumerate}
        \item
              \[
                  \sum\limits_{n=1}^m A_n^p
                  <\dfrac{p}{p-1}\sum\limits_{n=1}^mA_n^{p-1}a_n,
                  \qquad p>1,
              \]
              其中 $\{a_n\}$ 单调递增,$\displaystyle A_n=\dfrac{a_1+a_2+\cdots+a_n}{n}$;
        \item (Hardy-Landan 不等式)设 $p>1$,$\{a_n\}$ 同 (1),则
              \[
                  \sum\limits_{n=1}^m\left(\dfrac{a_1+a_2+\cdots+a_n}{n}\right)^p
                  <\left(\dfrac{p}{p-1}\right)^p\sum\limits_{n=1}^ma_n^p;
              \]
        \item (Cauleman 不等式)设 $\displaystyle\sum\limits_{n=1}^{\infty}a_n$ 是正项收敛级数,则
              \[
                  \sum\limits_{n=1}^{\infty}\sqrt[n]{a_1a_2\cdots a_n}
                  \leq e\sum\limits_{n=1}^{\infty}a_n.
              \]
    \end{enumerate}
\end{example}
\exampleblank

\begin{solution}
    \begin{enumerate}
        \item 由 $nA_n=(n-1)A_{n-1}+a_n$ 及凸函数 $x^p$ 的切线不等式,
              \[
                  nA_n^p-(n-1)A_{n-1}^p
                  \leq pA_n^{p-1}a_n-(p-1)A_n^p.
              \]
              对 $n=1,\ldots,m$ 求和并移项,即得题中不等式.
        \item 对 (1) 的右端使用 H\"older 不等式:
              \[
                  \sum\limits_{n=1}^{m}A_n^{p-1}a_n
                  \leq
                  \left(\sum\limits_{n=1}^{m}A_n^p\right)^{\frac{p-1}{p}}
                  \left(\sum\limits_{n=1}^{m}a_n^p\right)^{\frac{1}{p}}.
              \]
              与 (1) 合并并约去公共因子,得到所述 Hardy--Landau 不等式.
        \item 在 (2) 中以 $a_n^{1/p}$ 代替 $a_n$,得到
              \[
                  \sum\limits_{n=1}^{m}
                  \left(\dfrac{a_1^{1/p}+\cdots+a_n^{1/p}}n\right)^p
                  \leq\left(\dfrac{p}{p-1}\right)^p\sum\limits_{n=1}^{m}a_n.
              \]
              几何平均不超过算术平均,故左端逐项不小于
              $(a_1\cdots a_n)^{1/n}$.先令 $m\to\infty$,再令 $p\to\infty$,并用
              $(p/(p-1))^p\to e$,即得 Carleman 不等式.
    \end{enumerate}
\end{solution}

%例题8.2.11
\begin{example}
    研究
    \[
        \sum\limits_{n=1}^{\infty}\dfrac{1}{x_n^2}
    \]
    的敛散性,这里 $x_n$ 是方程 $x=\tan x$ 的正根,并且按递增的顺序编号.
\end{example}
\exampleblank

\begin{solution}
    第 $n$ 个正根位于 $(n\pi,n\pi+\pi/2)$ 内,因而 $x_n>n\pi$.于是
    \[
        0<\dfrac{1}{x_n^2}<\dfrac{1}{\pi^2n^2}.
    \]
    与 $p$ 级数比较可知原级数收敛.
\end{solution}

%例题8.2.12
\begin{example}
    设 $\displaystyle\sum\limits_{n=1}^{\infty}a_n$ 为收敛的正项级数,求证:
    \[
        \sum\limits_{n=1}^{\infty}a_n^{1-\frac{1}{n}}
    \]
    收敛.
\end{example}
\exampleblank

\begin{solution}
    因 $a_n\to0$,略去有限项后可设 $0<a_n<1$.分为两类:若
    $a_n\geq2^{-n}$,则 $a_n^{-1/n}\leq2$,故
    \[
        a_n^{1-\frac{1}{n}}\leq2a_n;
    \]
    若 $a_n<2^{-n}$,则
    \[
        a_n^{1-\frac{1}{n}}<2^{1-n}.
    \]
    所以
    \[
        a_n^{1-\frac{1}{n}}\leq2a_n+2^{1-n},
    \]
    比较可知题中级数收敛.
\end{solution}

%例题8.2.13
\begin{example}
    已知 $\varphi(x)$ 是 $(-\infty,+\infty)$ 上周期为 $1$ 的连续函数,满足 $\displaystyle\int_0^1\varphi(x)\,\mathrm{d}x=0$.设
    \[
        a_n=\int_0^1e^x\varphi(nx)\,\mathrm{d}x,
    \]
    求证:级数 $\displaystyle\sum\limits_{n=1}^{\infty}a_n^2$ 收敛.
\end{example}
\exampleblank

\begin{solution}
    令
    \[
        \Phi(t)=\int_0^t\varphi(u)\,\mathrm{d}u.
    \]
    由周期性及一个周期上的积分为零,$\Phi$ 是有界周期函数.换元 $t=nx$ 并分部积分,得
    \[
        \begin{aligned}
            a_n & =\dfrac{1}{n}\int_0^ne^{t/n}\varphi(t)\,\mathrm{d}t \\
                & =\dfrac{1}{n}\left[e^{t/n}\Phi(t)\right]_0^n
            -\dfrac{1}{n^2}\int_0^ne^{t/n}\Phi(t)\,\mathrm{d}t.
        \end{aligned}
    \]
    因 $\Phi(n)=\Phi(0)=0$,故 $|a_n|\leq C/n$.于是
    $\displaystyle \sum a_n^2$ 由 $\displaystyle \sum C^2/n^2$ 控制而收敛.
\end{solution}

%例题8.2.14
\begin{example}[\markstar]
    设
    \[
        \xi(s)=\sum\limits_{n=1}^{\infty}\dfrac{1}{n^s},
        \qquad s>1,
    \]
    则
    \[
        \xi(s)=s\int_1^{+\infty}\dfrac{[x]}{x^{s+1}}\,\mathrm{d}x.
    \]
\end{example}
\exampleblank

\begin{solution}
    在每个区间 $[n,n+1)$ 上有 $[x]=n$,故
    \[
        \begin{aligned}
            s\int_1^{+\infty}\dfrac{[x]}{x^{s+1}}\,\mathrm{d}x
             & =\sum\limits_{n=1}^{\infty}sn\int_n^{n+1}x^{-s-1}\,\mathrm{d}x             \\
             & =\sum\limits_{n=1}^{\infty}n\left(\dfrac{1}{n^s}-\dfrac{1}{(n+1)^s}\right) \\
             & =\sum\limits_{n=1}^{\infty}\dfrac{1}{n^s}=\xi(s),
        \end{aligned}
    \]
    其中最后一步由部分和裂项后令上限趋于无穷得到.
\end{solution}

%例题8.2.15
\begin{example}[\markstar]
    \begin{enumerate}
        \item 求证:当 $s>0$ 时,$\displaystyle\int_1^{+\infty}\dfrac{x-[x]}{x^{s+1}}\,\mathrm{d}x$ 收敛;
        \item 求证:当 $s>1$ 时,
              \[
                  \int_1^{+\infty}\dfrac{x-[x]}{x^{s+1}}\,\mathrm{d}x
                  =\dfrac{1}{s-1}-\dfrac{1}{s}\sum\limits_{n=1}^{\infty}\dfrac{1}{n^s}.
              \]
    \end{enumerate}
\end{example}
\exampleblank

\begin{solution}
    \begin{enumerate}
        \item 因 $0\leq x-[x]<1$,被积函数绝对值不超过 $x^{-s-1}$,故 $s>0$ 时绝对收敛.
        \item 当 $s>1$ 时,利用上一题的恒等式,
              \[
                  \begin{aligned}
                      \int_1^{+\infty}\dfrac{x-[x]}{x^{s+1}}\,\mathrm{d}x
                       & =\int_1^{+\infty}x^{-s}\,\mathrm{d}x
                      -\int_1^{+\infty}\dfrac{[x]}{x^{s+1}}\,\mathrm{d}x                       \\
                       & =\dfrac{1}{s-1}-\dfrac{1}{s}\sum\limits_{n=1}^{\infty}\dfrac{1}{n^s}.
                  \end{aligned}
              \]
    \end{enumerate}
\end{solution}

%例题8.2.16
\begin{example}
    设 $k>0,a>0$,证明:
    \begin{enumerate}
        \item $\displaystyle\int_a^{+\infty}\dfrac{\sin2n\pi x}{x^k}\,\mathrm{d}x$ 收敛;
        \item $\displaystyle\sum\limits_{n=1}^{\infty}\dfrac{1}{n}\dfrac{\sin2n\pi x}{x^k}$ 收敛.
    \end{enumerate}
\end{example}
\exampleblank

\begin{solution}
    \begin{enumerate}
        \item 对任意固定正整数 $n$,分部积分给出
              \[
                  \begin{aligned}
                      \int_a^R\dfrac{\sin(2n\pi x)}{x^k}\,\mathrm{d}x
                       & =\left[-\dfrac{\cos(2n\pi x)}{2n\pi x^k}\right]_a^R
                      -\dfrac{k}{2n\pi}\int_a^R\dfrac{\cos(2n\pi x)}{x^{k+1}}\,\mathrm{d}x.
                  \end{aligned}
              \]
              右端在 $R\to\infty$ 时有极限,故积分收敛.
        \item 记上述反常积分为 $I_n$.同一公式给出 $|I_n|\leq C/n$,其中 $C$ 与 $n$ 无关.因此
              \[
                  \sum\limits_{n=1}^{\infty}\dfrac{|I_n|}{n}
                  \leq C\sum\limits_{n=1}^{\infty}\dfrac{1}{n^2}<\infty.
              \]
              原题第二式把积分变量 $x$ 留在了求和号内;按``先取该反常积分再对 $n$ 求和''的自然理解,结论如上.
    \end{enumerate}
\end{solution}

%例题8.2.17
\begin{example}
    判别
    \[
        \sum\limits_{k=1}^{\infty}\left[e-\left(1+\dfrac{1}{1!}+\dfrac{1}{2!}+\cdots+\dfrac{1}{n!}\right)\right]
    \]
    的敛散性.
\end{example}
\exampleblank

\begin{solution}
    题面求和指标应为 $n$.记
    \[
        r_n=e-\sum\limits_{k=0}^{n}\dfrac{1}{k!}
        =\sum\limits_{k=n+1}^{\infty}\dfrac{1}{k!}.
    \]
    当 $n\geq1$ 时,
    \[
        0<r_n\leq\dfrac{1}{(n+1)!}
        \sum\limits_{j=0}^{\infty}\dfrac{1}{(n+2)^j}
        \leq\dfrac{2}{(n+1)!}.
    \]
    故 $\displaystyle \sum r_n$ 与阶乘级数比较可知收敛.
\end{solution}



\methodsection{(B) 通项趋势比较法}

\begin{theorem}
    设 $\displaystyle(A)\sim\sum\limits_{n=1}^{\infty}a_n$,$\displaystyle(B)\sim\sum\limits_{n=1}^{\infty}b_n$ 是两个正项级数,若存在极限
    \[
        \lim\limits_{n\to\infty}\dfrac{a_n}{b_n}=l,
    \]
    则当 $l>0$ 时,级数 $(A)$ 与 $(B)$ 同敛散.
\end{theorem}

%例题8.2.18
\begin{example}
    判别级数 $\displaystyle I=\sum\limits_{n=1}^{\infty}a_n$ 的敛散性:
    \begin{enumerate}
        \item $\displaystyle a_n=\dfrac{1}{n^{1+\frac{1}{n}}}$;
        \item $\displaystyle a_n=\dfrac{\sqrt[n]{2}-1}{n^p}$;
        \item $\displaystyle a_n=\dfrac{\sqrt{n+1}-\sqrt n}{n^p}$;
        \item $\displaystyle a_n=\left(1-\sqrt[3]{\dfrac{n-1}{n+1}}\right)^p\quad(p>0)$;
        \item $\displaystyle a_n=\left[e-\left(1+\dfrac{1}{n}\right)^n\right]^p$.
    \end{enumerate}
\end{example}
\exampleblank

\begin{solution}
    逐项作等价比较:
    \[
        \begin{gathered}
            n^{-1-\frac{1}{n}}\sim\dfrac{1}{n},
            \qquad
            \dfrac{\sqrt[n]2-1}{n^p}\sim\dfrac{\ln2}{n^{p+1}},\\
            \dfrac{\sqrt{n+1}-\sqrt n}{n^p}\sim\dfrac{1}{2n^{p+\frac{1}{2}}},
            \qquad
            \left(1-\sqrt[3]{\dfrac{n-1}{n+1}}\right)^p
            \sim\left(\dfrac{2}{3n}\right)^p,\\
            \left[e-\left(1+\dfrac{1}{n}\right)^n\right]^p
            \sim\left(\dfrac{e}{2n}\right)^p.
        \end{gathered}
    \]
    所以 (1) 发散;(2) 在 $p>0$ 时收敛,$p\leq0$ 时发散;(3) 在
    $p>1/2$ 时收敛,否则发散;(4),(5) 均在 $p>1$ 时收敛,在
    $0<p\leq1$ 时发散.
\end{solution}

%例题8.2.19
\begin{example}
    判断级数 $\displaystyle I=\sum\limits_{n=1}^{\infty}a_n$ 的敛散性:
    \begin{enumerate}
        \item $\displaystyle a_n=\left(\dfrac{1}{n}-\sin\dfrac{1}{n}\right)^p\quad(p>0)$;
        \item $\displaystyle a_n=\left(e^{\frac{1}{n}}-\sin\dfrac{1}{n}\right)^{np}-1$;
        \item $\displaystyle a_n=\ln\left(\dfrac{1}{\cos\dfrac{2\pi}{n}}\right)\quad(n>4)$.
    \end{enumerate}
\end{example}
\exampleblank

\begin{solution}
    Taylor 展开给出
    \[
        \dfrac{1}{n}-\sin\dfrac{1}{n}\sim\dfrac{1}{6n^3},
        \qquad
        -\ln\cos\dfrac{2\pi}{n}\sim\dfrac{2\pi^2}{n^2}.
    \]
    故 (1) 在 $p>1/3$ 时收敛,在 $0<p\leq1/3$ 时发散;(3) 收敛.

    对 (2),
    \[
        e^{1/n}-\sin\dfrac{1}{n}=1+\dfrac{1}{2n^2}+O\left(\dfrac{1}{n^3}\right),
    \]
    从而对固定 $p\neq0$,
    \[
        \left(e^{1/n}-\sin\dfrac{1}{n}\right)^{np}-1
        \sim\dfrac{p}{2n}.
    \]
    因此该级数发散;$p=0$ 时各项为零.
\end{solution}

%例题8.2.20
\begin{example}
    设正数列 $\{a_n\}$ 满足
    \[
        \lim\limits_{n\to\infty}a_n\sum\limits_{i=1}^na_i^2=1.
    \]
    试证明:
    \[
        \lim\limits_{n\to\infty}\sqrt[3]{3na_n}=1,
    \]
    且 $\displaystyle\sum\limits_{n=1}^{\infty}a_n^3$ 发散.
\end{example}
\exampleblank

\begin{solution}
    记 $\displaystyle B_n=\sum_{i=1}^{n}a_i^2$.题设等价于 $a_nB_n\to1$,故
    $B_n\to\infty$ 且 $a_n\sim1/B_n$.于是
    \[
        \begin{aligned}
            B_n^3-B_{n-1}^3
             & =(B_n-B_{n-1})(B_n^2+B_nB_{n-1}+B_{n-1}^2)          \\
             & =a_n^2(B_n^2+B_nB_{n-1}+B_{n-1}^2)\longrightarrow3.
        \end{aligned}
    \]
    由 Stolz 定理,$B_n^3/n\to3$,进而
    \[
        a_n\sim\dfrac{1}{\sqrt[3]{3n}},
        \qquad a_n^3\sim\dfrac{1}{3n}.
    \]
    所以 $\displaystyle \sum a_n^3$ 发散.与此同时
    \[
        \sqrt[3]{3n}\,a_n\longrightarrow1,
    \]
    而题面写成 $\sqrt[3]{3na_n}\to1$,后者与上述必然渐近式矛盾,故根号范围应核对原始资料.
\end{solution}



\subsection{特殊比较判别}


\methodsection{(A) Cauchy 凝聚判别法}

\begin{theorem}[\markstar\ Cauchy 凝聚判别法]
    设 $\{a_n\}$ 是递减正数列,则
    \[
        \sum\limits_{n=1}^{\infty}a_n\text{ 收敛}
        \quad\Longleftrightarrow\quad
        \sum\limits_{n=0}^{\infty}2^na_{2^n}\text{ 收敛}.
    \]
\end{theorem}

\begin{proposition}[\markstar]
    设 $\{a_n\}$ 是递减正数列,若 $\displaystyle\sum\limits_{n=1}^{\infty}a_n$ 收敛,则
    \[
        \lim\limits_{n\to\infty}na_n=0.
    \]
\end{proposition}

\begin{remark}
    反常积分亦有类似结论,甚至证明方法都完全相似.
\end{remark}

%例题8.2.21
\begin{example}
    设 $\displaystyle\sum\limits_{n=1}^{\infty}a_n$ 是正项收敛级数,若 $\{na_n\}$ 为单调数列,则
    \[
        \lim\limits_{n\to\infty}na_n\ln n=0.
    \]
\end{example}
\exampleblank

\begin{solution}
    置 $b_n=na_n$.若 $b_n$ 单调递增,则因其为正数,存在 $c>0$ 使最终
    $a_n=b_n/n\geq c/n$,与级数收敛矛盾.因此 $b_n$ 必单调递减.
    对充分大的 $n$,
    \[
        \sum\limits_{k=\lceil\sqrt n\rceil}^{n}a_k
        =\sum\limits_{k=\lceil\sqrt n\rceil}^{n}\dfrac{b_k}{k}
        \geq b_n\sum\limits_{k=\lceil\sqrt n\rceil}^{n}\dfrac{1}{k}.
    \]
    左端作为收敛级数的尾和趋于零,而右侧调和和渐近于 $\frac{1}{2}\ln n$,故
    $b_n\ln n=na_n\ln n\to0$.
\end{solution}

%例题8.2.22
\begin{example}
    判别下列级数的敛散性:
    \begin{enumerate}
        \item $\displaystyle\sum\limits_{n=1}^{\infty}\dfrac{1}{n^p}$;
        \item $\displaystyle\sum\limits_{n=2}^{\infty}\dfrac{1}{n(\ln n)^p}$;
        \item $\displaystyle\sum\limits_{n=3}^{\infty}\dfrac{1}{n\ln n\ln\ln n}$.
    \end{enumerate}
\end{example}
\exampleblank

\begin{solution}
    由 Cauchy 凝聚判别:
    \begin{enumerate}
        \item $\displaystyle \sum n^{-p}$ 与 $\displaystyle \sum2^{n(1-p)}$ 同敛散,故恰在 $p>1$ 时收敛.
        \item $\displaystyle \sum1/[n(\ln n)^p]$ 与常数倍的 $\displaystyle \sum n^{-p}$ 同敛散,故恰在 $p>1$ 时收敛.
        \item 再作一次凝聚,或直接与积分比较,可化为调和型级数
              $\displaystyle \sum1/(n\ln n)$,故发散.
    \end{enumerate}
\end{solution}

%例题8.2.23
\begin{example}
    试证明下列命题:
    \begin{enumerate}
        \item 设 $\{a_n\}$ 是递减正数列,若 $\displaystyle\sum\limits_{n=1}^{\infty}a_{2^n}$ 发散,则 $\displaystyle\sum\limits_{n=1}^{\infty}\dfrac{a_n}{n}$ 发散;
        \item 设 $\{a_n\}$ 是递减于零的数列:
              \begin{enumerate}[label=(\roman*)]
                  \item 若 $\displaystyle\sum\limits_{n=1}^{\infty}a_n=+\infty$,则 $\displaystyle\sum\limits_{n=1}^{\infty}\min\left\{a_n,\dfrac{1}{n}\right\}=+\infty$;
                  \item 若 $\displaystyle\sum\limits_{n=1}^{\infty}\dfrac{a_n}{n}=+\infty$,则 $\displaystyle\sum\limits_{n=1}^{\infty}\dfrac{1}{n}\min\left\{a_n,\dfrac{1}{\ln n}\right\}=+\infty$.
              \end{enumerate}
    \end{enumerate}
\end{example}
\exampleblank

\begin{solution}
    \begin{enumerate}
        \item 在区间 $2^k\leq n<2^{k+1}$ 上,
              \[
                  \sum\limits_{n=2^k}^{2^{k+1}-1}\dfrac{a_n}{n}
                  \geq\dfrac{1}{2}a_{2^{k+1}}.
              \]
              右侧分块和发散,故原级数发散.
        \item 令 $b_n=\min\{a_n,1/n\}$.若 $\displaystyle \sum b_n$ 收敛,则因 $b_n$ 递减,有 $nb_n\to0$.于是充分大时不可能有 $b_n=1/n$,故最终 $b_n=a_n$,这与 $\displaystyle \sum a_n$ 发散矛盾.

              再令 $c_n=\min\{a_n,1/\ln n\}$.若 $\displaystyle \sum c_n/n$ 收敛,则对递减数列 $c_n$ 使用凝聚判别可得 $c_n\ln n\to0$.故最终 $c_n<1/\ln n$,即 $c_n=a_n$,与 $\displaystyle \sum a_n/n$ 发散矛盾.
    \end{enumerate}
\end{solution}

%例题8.2.24
\begin{example}[\markstar\ 推广的 Cauchy 凝聚定理]
    设 $\{a_n\}$ 是严格递减的正数列,若存在严格递增的正整数列 $\{m_k\}$ 和 $M>0$,使得
    \[
        m_{k+1}-m_k\leq M(m_k-m_{k-1}),
        \qquad k=2,3,\ldots,
    \]
    则级数
    \[
        \sum\limits_{n=1}^{\infty}a_n
        \quad\text{与}\quad
        \sum\limits_{k=1}^{\infty}(m_{k+1}-m_k)a_{m_k}
    \]
    同敛散.
\end{example}
\exampleblank

\begin{solution}
    记 $d_k=m_{k+1}-m_k$.利用 $a_n$ 的单调性,
    \[
        d_ka_{m_{k+1}}
        \leq\sum\limits_{n=m_k}^{m_{k+1}-1}a_n
        \leq d_ka_{m_k}.
    \]
    因此 $\displaystyle \sum d_ka_{m_k}$ 收敛立即推出 $\displaystyle \sum a_n$ 收敛.反之,将左端下界改写为
    \[
        d_ka_{m_{k+1}}
        \geq\dfrac{1}{M} d_{k+1}a_{m_{k+1}},
    \]
    可见 $\displaystyle \sum a_n$ 收敛蕴含 $\displaystyle \sum d_{k+1}a_{m_{k+1}}$ 收敛.移去有限个首项不影响敛散性,故两级数同敛散.
\end{solution}

\begin{remark}
    因此若 $\{a_n\}$ 是递减正数列,则 $\displaystyle\sum\limits_{n=1}^{\infty}a_n$ 与下列之一同敛散:
    \[
        \sum\limits_{n=1}^{\infty}3^na_{3^n},
        \qquad
        \sum\limits_{n=1}^{\infty}na_{n^2},
        \qquad
        \sum\limits_{n=1}^{\infty}n^2a_{n^3},
        \qquad
        \sum\limits_{n=1}^{\infty}e^na_{e^n}.
    \]
\end{remark}

%例题8.2.25
\begin{example}
    试证明下列命题:
    \begin{enumerate}
        \item 设 $\displaystyle\sum\limits_{n=1}^{\infty}a_n$ 是收敛正项级数,记 $\displaystyle R_n=\sum\limits_{k=n+1}^{\infty}a_k$,若 $\displaystyle\sum\limits_{n=1}^{\infty}R_n$ 收敛,则 $\displaystyle\lim\limits_{n\to\infty}na_n=0$;
        \item 设 $\{a_n\}$ 是递减正数列,若 $\displaystyle\sum\limits_{n=1}^{\infty}\dfrac{a_n}{\sqrt n}$ 收敛,则级数 $\displaystyle\sum\limits_{n=1}^{\infty}a_n^2$ 收敛.
    \end{enumerate}
\end{example}
\exampleblank

\begin{solution}
    \begin{enumerate}
        \item 交换非负项的求和次序,
              \[
                  \sum\limits_{n=1}^{\infty}R_n
                  =\sum\limits_{n=1}^{\infty}\sum\limits_{k=n+1}^{\infty}a_k
                  =\sum\limits_{k=2}^{\infty}(k-1)a_k<\infty.
              \]
              因而正项级数 $\displaystyle \sum(k-1)a_k$ 的通项趋于零,即 $na_n\to0$.
        \item 由单调性,
              \[
                  \sum\limits_{k=\lceil n/2\rceil}^{n}\dfrac{a_k}{\sqrt k}
                  \geq\dfrac{n}{2}\dfrac{a_n}{\sqrt n}=\dfrac{1}{2}\sqrt n,a_n.
              \]
              左端是收敛级数的尾和,故 $\sqrt n,a_n\to0$.最终有
              $a_n\leq1/\sqrt n$,于是 $a_n^2\leq a_n/\sqrt n$,比较即得收敛.
    \end{enumerate}
\end{solution}

%例题8.2.26
\begin{example}
    设 $\{a_n\}$ 递减于零,若 $b_n=a_n-2a_{n+1}+a_{n+2}\geq0$($n\in\mathbb{N}$),则
    \[
        \sum\limits_{n=1}^{\infty}nb_n=a_1.
    \]
\end{example}
\exampleblank

\begin{solution}
    置 $d_n=a_n-a_{n+1}\geq0$,则 $b_n=d_n-d_{n+1}\geq0$,所以
    $d_n$ 单调递减.有限求和得
    \[
        \begin{aligned}
            \sum\limits_{n=1}^{N}nb_n
             & =\sum\limits_{n=1}^{N}n(d_n-d_{n+1}) \\
             & =\sum\limits_{n=1}^{N}d_n-Nd_{N+1}   \\
             & =a_1-a_{N+1}-Nd_{N+1}.
        \end{aligned}
    \]
    又 $\displaystyle \sum d_n=a_1<\infty$ 且 $d_n$ 递减,故 $Nd_{N+1}\to0$;同时
    $a_{N+1}\to0$.令 $N\to\infty$ 即得所求等式.
\end{solution}



\methodsection{(B) 比式判别法}

%例题8.2.27
\begin{example}
    讨论级数 $\displaystyle I=\sum\limits_{n=1}^{\infty}a_n$ 的敛散性:
    \begin{enumerate}
        \item $\displaystyle a_n=\dfrac{3^n n!}{n^n}$;
        \item $\displaystyle a_n=\dfrac{(2n-1)!!}{3^n n!}$;
        \item $\displaystyle a_n=\dfrac{(n!)^3}{3^{n^4/3}}$;
        \item $\displaystyle a_n=\dfrac{(\alpha+1)(2\alpha+1)\cdots(n\alpha+1)}{(\beta+1)(2\beta+1)\cdots(n\beta+1)}\quad(\alpha>0,\beta>0)$;
        \item $\displaystyle a_n=\dfrac{n^\alpha}{\beta^n}\quad(\beta>0)$.
    \end{enumerate}
\end{example}
\exampleblank

\begin{solution}
    逐项使用比式判别:
    \begin{enumerate}
        \item $a_{n+1}/a_n=3(n/(n+1))^n\to3/e>1$,故发散.
        \item $a_{n+1}/a_n=(2n+1)/(3n+3)\to2/3$,故收敛.
        \item $a_{n+1}/a_n\to0$,故收敛.
        \item 比值趋于 $\alpha/\beta$;当 $\alpha<\beta$ 时收敛,当
              $\alpha\geq\beta$ 时通项不趋于零或比值大于一,故发散.
        \item 当 $\beta>1$ 时收敛;当 $0<\beta<1$ 时通项不趋于零;当
              $\beta=1$ 时化为 $\displaystyle \sum n^\alpha$,恰在 $\alpha<-1$ 时收敛.
    \end{enumerate}
\end{solution}

%例题8.2.28
\begin{example}
    设 $a_1\geq a_2\geq\cdots\geq a_n\geq\cdots>0$,$\displaystyle\lim\limits_{n\to\infty}\dfrac{a_{2n}}{a_n}=\lambda$,则级数 $\displaystyle I=\sum\limits_{n=1}^{\infty}a_n$ 在 $\lambda<\dfrac{1}{2}$ 时收敛,在 $\lambda>\dfrac{1}{2}$ 时发散.
\end{example}
\exampleblank

\begin{solution}
    由 Cauchy 凝聚判别,$\displaystyle \sum a_n$ 与 $\displaystyle \sum2^ka_{2^k}$ 同敛散,而
    \[
        \dfrac{2^{k+1}a_{2^{k+1}}}{2^ka_{2^k}}
        =2\dfrac{a_{2^{k+1}}}{a_{2^k}}\longrightarrow2\lambda.
    \]
    故 $\lambda<1/2$ 时凝聚级数收敛,$\lambda>1/2$ 时发散.
\end{solution}

%例题8.2.29
\begin{example}
    \begin{enumerate}
        \item 设 $a_n>0$,$\displaystyle\lim\limits_{n\to\infty}\dfrac{a_{2n}}{a_{2n-1}}=\alpha$,$\displaystyle\lim\limits_{n\to\infty}\dfrac{a_{2n+1}}{a_{2n}}=\beta$,若 $\alpha\beta<1$,则 $\displaystyle I=\sum\limits_{n=1}^{\infty}a_n$ 收敛;
        \item 设 $a_n>0$,$\displaystyle\lim\limits_{n\to\infty}\dfrac{a_{n+1}}{a_n}=\lambda<1$,则 $\displaystyle\lim\limits_{n\to\infty}\dfrac{a_{2n}}{a_n}=0$.
    \end{enumerate}
\end{example}
\exampleblank

\begin{solution}
    \begin{enumerate}
        \item 有
              \[
                  \dfrac{a_{2n+1}}{a_{2n-1}}
                  =\dfrac{a_{2n+1}}{a_{2n}}\dfrac{a_{2n}}{a_{2n-1}}
                  \longrightarrow\alpha\beta<1.
              \]
              故奇数项子级数由比式判别收敛;又 $a_{2n}/a_{2n-1}\to\alpha$,偶数项子级数也收敛.
        \item 任取 $q\in(\lambda,1)$,充分大时 $a_{k+1}/a_k\leq q$,从而
              \[
                  0<\dfrac{a_{2n}}{a_n}=\prod\limits_{k=n}^{2n-1}\dfrac{a_{k+1}}{a_k}
                  \leq q^n\longrightarrow0.
              \]
    \end{enumerate}
\end{solution}

%例题8.2.30
\begin{example}[\markstar\ 比式判别法推广]
    设 $a_n>0$,$k_0$ 是固定正整数,若
    \[
        \lim\limits_{n\to\infty}\dfrac{a_{n+k_0}}{a_n}=\lambda,
    \]
    则 $\displaystyle I=\sum\limits_{n=1}^{\infty}a_n$ 在 $\lambda<1$ 时收敛,在 $\lambda>1$ 时发散.
\end{example}
\exampleblank

\begin{solution}
    将级数按模 $k_0$ 分为 $k_0$ 个子级数.对每个余数类,其相邻项之比均趋于 $\lambda$.若 $\lambda<1$,每个子级数均由比式判别收敛;若 $\lambda>1$,每个余数类上的项最终递增,原通项不能趋于零,故级数发散.
\end{solution}

%例题8.2.31
\begin{example}
    设 $\displaystyle I=\sum\limits_{n=1}^{\infty}a_n$ 是正项级数,若
    \[
        \lim\limits_{n\to\infty}\left(\dfrac{a_{n+1}}{a_n}\right)^n=\lambda,
    \]
    则 $I$ 在 $\lambda<e^{-1}$ 时收敛,$\lambda>e^{-1}$ 时发散.
\end{example}
\exampleblank

\begin{solution}
    取对数得
    \[
        n\ln\dfrac{a_{n+1}}{a_n}\longrightarrow\ln\lambda.
    \]
    于是
    \[
        n\left(\dfrac{a_n}{a_{n+1}}-1\right)
        \longrightarrow-\ln\lambda.
    \]
    应用 Raabe 判别:若 $-\ln\lambda>1$,即 $\lambda<e^{-1}$,则收敛;若
    $-\ln\lambda<1$,即 $\lambda>e^{-1}$,则发散.
\end{solution}

%例题8.2.32
\begin{example}
    判别下列级数的敛散性:
    \begin{enumerate}
        \item $\displaystyle\dfrac{1}{1+a}+\dfrac{2}{1+a^2}+\dfrac{4}{1+a^4}+\dfrac{8}{1+a^8}+\cdots$;
        \item $\displaystyle\dfrac{a}{1+a}+\dfrac{2a^2}{1+a^2}+\dfrac{4a^4}{1+a^4}+\cdots$.
    \end{enumerate}
\end{example}
\exampleblank

\begin{solution}
    \begin{enumerate}
        \item 通项为 $2^{n-1}/(1+a^{2^{n-1}})$.当 $a>1$ 时其根式判别上极限为 $1/a<1$,故收敛;当 $0<a\leq1$ 时通项不趋于零.故恰在 $a>1$ 时收敛.
        \item 通项为 $2^{n-1}a^{2^{n-1}}/(1+a^{2^{n-1}})$.当 $0<a<1$ 时根式判别给出收敛;当 $a\geq1$ 时通项不趋于零.故恰在 $0<a<1$ 时收敛.
    \end{enumerate}
\end{solution}



\methodsection{(C) 根式判别法}

%例题8.2.33
\begin{example}
    判别级数 $\displaystyle I=\sum\limits_{n=1}^{\infty}a_n$ 的敛散性:
    \begin{enumerate}
        \item $\displaystyle a_n=(\sqrt[n]{n}-1)^n$;
        \item $\displaystyle a_n=\left(\dfrac{n}{n+1}\right)^{n(n+1)}$;
        \item $\displaystyle a_n=\left(\dfrac{1+\cos n}{2+\cos n}\right)^{2n-\ln n}$;
        \item $\displaystyle a_n=\dfrac{n^3}{(\sqrt2+(-1)^n)^n3^n}$;
        \item $\displaystyle a_n=\dfrac{\ln(1+\alpha^n)}{n^\beta}\quad(\alpha\geq0)$.
    \end{enumerate}
\end{example}
\exampleblank

\begin{solution}
    使用根式判别:
    \begin{enumerate}
        \item $\sqrt[n]{a_n}=\sqrt[n]n-1\sim(\ln n)/n\to0$,故收敛.
        \item $\sqrt[n]{a_n}=(n/(n+1))^{n+1}\to e^{-1}$,故收敛.
        \item
              \[
                  \limsup\sqrt[n]{a_n}
                  \leq\max_{-1\leq t\leq1}\left(\dfrac{1+t}{2+t}\right)^2
                  =\dfrac{4}{9}<1,
              \]
              故收敛.
        \item 奇偶子列的 $n$ 次方根分别趋于
              $1/[3(\sqrt2+1)]$ 与 $1/[3(\sqrt2-1)]$,二者均小于 $1$,故收敛.
        \item 当 $0\leq\alpha<1$ 时,通项至多为常数倍的
              $\alpha^n/n^\beta$,故收敛;当 $\alpha=1$ 时,恰在 $\beta>1$ 时收敛;当 $\alpha>1$ 时,
              \[
                  \ln(1+\alpha^n)\sim n\ln\alpha,
              \]
              故恰在 $\beta>2$ 时收敛.
    \end{enumerate}
\end{solution}

%例题8.2.34
\begin{example}
    设 $\{a_n\},\{b_n\}$ 为正数列,若有
    \[
        \lim\limits_{n\to\infty}\dfrac{a_{n+1}}{a_n}=\alpha,
        \qquad
        \lim\limits_{n\to\infty}\sqrt[n]{b_n}=\beta,
    \]
    试证明级数 $\displaystyle I=\sum\limits_{n=1}^{\infty}a_nb_n$ 在 $\alpha\beta<1$ 时收敛,$\alpha\beta>1$ 时发散.
\end{example}
\exampleblank

\begin{solution}
    由 $a_{n+1}/a_n\to\alpha$ 可得 $\sqrt[n]{a_n}\to\alpha$.于是
    \[
        \sqrt[n]{a_nb_n}=\sqrt[n]{a_n}\sqrt[n]{b_n}\longrightarrow\alpha\beta.
    \]
    结论立即由根式判别得到.
\end{solution}



\methodsection{(D) Cauchy 积分判别法}

%例题8.2.35
\begin{example}
    证明下列命题:
    \begin{enumerate}
        \item 设 $a_1>1$,$\displaystyle a_{n+1}=\dfrac{a_n}{1+a_n^\alpha}$,$0<\alpha<1$,则 $\displaystyle\sum\limits_{n=1}^{\infty}a_n$ 收敛;
        \item 设 $a_n>0$,$\displaystyle\sum\limits_{n=1}^{\infty}\dfrac{a_n}{n}$ 收敛,则
              \[
                  \sum\limits_{n=1}^{\infty}\left(\sum\limits_{k=1}^{\infty}\dfrac{a_n}{n^2+k^2}\right)
              \]
              收敛.
    \end{enumerate}
\end{example}
\exampleblank

\begin{solution}
    \begin{enumerate}
        \item 数列 $a_n$ 递减至零.令 $u_n=a_n^{-\alpha}$,则
              \[
                  u_{n+1}-u_n
                  =a_n^{-\alpha}\left[(1+a_n^\alpha)^\alpha-1\right]\longrightarrow\alpha.
              \]
              由 Stolz 定理,$u_n/n\to\alpha$,故
              $a_n\sim(\alpha n)^{-1/\alpha}$.因 $0<\alpha<1$,有 $1/\alpha>1$,所以 $\displaystyle \sum a_n$ 收敛.
        \item 当前内层求和的分子为 $a_n$.利用积分比较,
              \[
                  \sum\limits_{k=1}^{\infty}\dfrac{1}{n^2+k^2}
                  \leq\int_0^{+\infty}\dfrac{\mathrm{d}x}{n^2+x^2}
                  =\dfrac{\pi}{2n}.
              \]
              因此题中双重级数不超过 $\displaystyle (\pi/2)\sum a_n/n$,从而收敛.
    \end{enumerate}
\end{solution}

%例题8.2.36
\begin{example}
    设
    \[
        a_n=2\sqrt n-\sum\limits_{k=1}^n\dfrac{1}{\sqrt k},
    \]
    则 $\displaystyle\lim\limits_{n\to\infty}a_n=l$ 且 $1\leq l\leq2$.
\end{example}
\exampleblank

\begin{solution}
    有
    \[
        a_{n+1}-a_n
        =\dfrac{2}{\sqrt{n+1}+\sqrt n}-\dfrac{1}{\sqrt{n+1}}>0,
    \]
    故 $a_n$ 单调递增.又由递减函数 $x^{-1/2}$ 的积分比较,
    \[
        2(\sqrt{n+1}-1)leq\sum\limits_{k=1}^{n}\dfrac{1}{\sqrt k}
        \leq1+2(\sqrt n-1).
    \]
    因此
    \[
        1\leq a_n<2.
    \]
    单调有界性给出 $a_n\to l$,且 $1\leq l\leq2$.
\end{solution}

%例题8.2.37
\begin{example}
    设 $f(x)$ 于 $\lbrack 1,+\infty)$ 上可导,$f'(x)$ 单调递增,且 $f(x)\to A$($x\to+\infty$),证明:
    \[
        \sum\limits_{n=1}^{\infty}f'(n)
    \]
    收敛.
\end{example}
\exampleblank

\begin{solution}
    先证 $f'(x)\leq0$.若某点 $f'(x_0)>0$,由 $f'$ 单调递增,之后恒有
    $f'(x)\geq f'(x_0)>0$,这与 $f(x)$ 有有限极限矛盾.因此 $-f'(x)$ 是非负单调递减函数.由积分比较,
    \[
        \sum\limits_{n=2}^{\infty}|f'(n)|
        \leq\int_1^{+\infty}|f'(x)|\,\mathrm{d}x
        =f(1)-A<\infty.
    \]
    所以 $\displaystyle \sum f'(n)$ 绝对收敛.
\end{solution}



\subsection{其他比较判别 *}

%例题8.2.38
\begin{example}[比值对数判别法]
    设有正项级数 $\displaystyle\sum\limits_{n=1}^{\infty}a_n$,且存在
    \[
        \lim\limits_{n\to\infty}\left(n\ln\dfrac{a_n}{a_{n+1}}\right)=l,
        \qquad -\infty\leq l\leq+\infty.
    \]
    若 $l>1$,则级数收敛;若 $l<1$,则级数发散.
\end{example}
\exampleblank

\begin{solution}
    若 $l>1$,取 $1<p<l$,则充分大时
    \[
        \ln\dfrac{a_n}{a_{n+1}}>\dfrac{p}{n}.
    \]
    求和后得到 $a_n\leq Cn^{-p}$,故级数收敛.若 $l<1$,取
    $l<p<1$,同理有 $a_n\geq cn^{-p}$,故级数发散.
\end{solution}

%例题8.2.39
\begin{example}[Rabbe 判别法]
    设有正项级数 $\displaystyle I=\sum\limits_{n=1}^{\infty}a_n$ 满足
    \[
        \lim\limits_{n\to\infty}n\left(\dfrac{a_n}{a_{n+1}}-1\right)=l,
        \qquad -\infty\leq l\leq+\infty.
    \]
    若 $l>1$,则 $I$ 收敛;若 $l<1$,则 $I$ 发散.
\end{example}
\exampleblank

\begin{solution}
    由 $u-1\sim\ln u$($u\to1$),当极限有限时,本判别与上一题的比值对数判别等价.更直接地,若 $l>1$,取 $1<p<l$,则充分大时
    \[
        \dfrac{a_n}{a_{n+1}}>1+\dfrac{p}{n}.
    \]
    与 $n^{-p}$ 作比式比较即得收敛;若 $l<1$,取 $l<p<1$ 作反向比较,得到发散.无穷极限情形同理.
\end{solution}

%例题8.2.40
\begin{example}[\markstar\ Kummer 判别法]
    设 $\displaystyle\sum\limits_{n=1}^{\infty}a_n,\sum\limits_{n=1}^{\infty}b_n$ 是正项级数,$\displaystyle\sum\limits_{n=1}^{\infty}b_n$ 发散,若有
    \[
        \lim\limits_{n\to\infty}\left(\dfrac{1}{b_n}\dfrac{a_n}{a_{n+1}}-\dfrac{1}{b_{n+1}}\right)=\lambda,
    \]
    则 $\displaystyle\sum\limits_{n=1}^{\infty}a_n$ 在 $\lambda>0$ 时收敛,在 $\lambda<0$ 时发散.
\end{example}
\exampleblank

\begin{solution}
    置 $q_n=1/b_n$.题设括号内的量为
    \[
        q_n\dfrac{a_n}{a_{n+1}}-q_{n+1}.
    \]
    若 $\lambda>0$,则充分大时它至少为某个 $c>0$,故
    \[
        q_na_n-q_{n+1}a_{n+1}\geq ca_{n+1}.
    \]
    求和后右端有界,从而 $\displaystyle \sum a_n$ 收敛.若 $\lambda<0$,则最终
    \[
        q_{n+1}a_{n+1}-q_na_n\geq ca_{n+1}>0.
    \]
    于是 $q_na_n$ 最终递增并有正下界,即 $a_n\geq Cb_n$;因
    $\displaystyle \sum b_n$ 发散,$\displaystyle \sum a_n$ 也发散.
\end{solution}

%例题8.2.41
\begin{example}
    设
    \[
        a_n=C_{\alpha}^{n}=\dfrac{\alpha(\alpha-1)\cdots(\alpha-n+1)}{n!}.
    \]
    \begin{enumerate}
        \item $\alpha$ 在什么条件下,$\displaystyle\lim\limits_{n\to\infty}a_n=0$;
        \item 讨论 $\displaystyle\sum\limits_{n=1}^{\infty}a_n$ 的敛散性.
    \end{enumerate}
\end{example}
\exampleblank

\begin{solution}
    若 $\alpha=m$ 为非负整数,则 $a_n=0$($n>m$).若 $\alpha$ 不是非负整数,利用 Gamma 函数的标准渐近式可得
    \[
        a_n=\dfrac{(-1)^n}{\Gamma(-\alpha)}n^{-\alpha-1}(1+o(1)).
    \]
    故 $a_n\to0$ 当且仅当 $\alpha>-1$.

    当 $\alpha>0$ 时级数绝对收敛;当 $-1<\alpha<0$ 时它是最终单调的交错级数,条件收敛;$\alpha=0$ 时各项均为零.由二项式展开的 Abel 极限定理,其和为 $2^\alpha-1$.因此题中级数恰在 $\alpha>-1$ 时收敛.
\end{solution}

%例题8.2.42
\begin{example}[对数判别法]
    设 $\displaystyle I=\sum\limits_{n=1}^{\infty}a_n$ 是正项级数,若存在
    \[
        \lim\limits_{n\to\infty}\dfrac{\ln\dfrac{1}{a_n}}{\ln n}=l,
    \]
    则当 $l>1$ 时,$I$ 收敛;当 $l<1$ 时,$I$ 发散.
\end{example}
\exampleblank

\begin{solution}
    若 $l>1$,取 $1<p<l$.充分大时
    \[
        \dfrac{\ln(1/a_n)}{\ln n}>p,
    \]
    即 $a_n<n^{-p}$,故收敛.若 $l<1$,取 $l<p<1$,则最终
    $a_n>n^{-p}$,故发散.
\end{solution}


\newpage

\section{一般级数}

\begin{proposition}[\marktwostar\ Abel 变换]
    设有 $\{a_n\},\{b_n\}$ 两个数列,$\displaystyle B_n=\sum\limits_{k=1}^nb_k$,则:
    \begin{enumerate}
        \item
              \[
                  \sum\limits_{k=1}^na_kb_k
                  =a_nB_n+\sum\limits_{k=1}^{n-1}(a_k-a_{k+1})B_k;
              \]
        \item
              \[
                  \sum\limits_{k=n+1}^ma_kb_k
                  =-A_nb_{n+1}+\sum\limits_{k=n+1}^{m-1}A_k(b_k-b_{k+1})+A_mb_m.
              \]
    \end{enumerate}
\end{proposition}

\begin{proposition}[\markstar]
    \begin{enumerate}
        \item
              \[
                  \left|\sum\limits_{k=1}^n\sin kx\right|
                  =\left|\sum\limits_{k=1}^n
                  \dfrac{\cos\dfrac{2k-1}{2}x-\cos\dfrac{2k+1}{2}x}{2\sin\dfrac{x}{2}}\right|
                  \leq\dfrac{1}{|\sin\dfrac{x}{2}|};
              \]
        \item
              \[
                  \left|\sum\limits_{k=1}^n\cos kx\right|
                  =\left|\sum\limits_{k=1}^n
                  \dfrac{\sin\dfrac{2k+1}{2}x-\sin\dfrac{2k-1}{2}x}{2\sin\dfrac{x}{2}}\right|
                  \leq\dfrac{1}{|\sin\dfrac{x}{2}|}.
              \]
    \end{enumerate}
\end{proposition}

%例题8.3.1
\begin{example}
    证明级数
    \[
        1-\dfrac{1}{3}\left(1+\dfrac{1}{2}\right)
        +\dfrac{1}{5}\left(1+\dfrac{1}{2}+\dfrac{1}{3}\right)
        -\dfrac{1}{7}\left(1+\dfrac{1}{2}+\dfrac{1}{3}+\dfrac{1}{4}\right)+\cdots
    \]
    是收敛的.
\end{example}
\exampleblank

\begin{solution}
    第 $n$ 项为
    \[
        (-1)^{n-1}\dfrac{H_n}{2n-1},
        \qquad H_n=\sum\limits_{k=1}^{n}\dfrac{1}{k}.
    \]
    因为 $H_n/(2n-1)\to0$,且该数列从某项起单调递减,所以由 Leibniz 判别法,原级数收敛.
\end{solution}

%例题8.3.2
\begin{example}[\markstar]
    \begin{enumerate}
        \item 证明 $\displaystyle\sum\limits_{n=1}^{\infty}\dfrac{(-1)^{[\sqrt n]}}{n}$ 收敛;
        \item 证明 $\displaystyle\sum\limits_{n=1}^{\infty}\dfrac{(-1)^{\sqrt[3]{n}}}{n}$ 收敛.
    \end{enumerate}
\end{example}
\exampleblank

\begin{solution}
    \begin{enumerate}
        \item 按 $k^2\leq n<(k+1)^2$ 分块.第 $k$ 块同号,且其和
              \[
                  B_k=(-1)^k\sum\limits_{n=k^2}^{(k+1)^2-1}\dfrac{1}{n}
                  =(-1)^k\left(\dfrac{2}{k}+O\left(\dfrac{1}{k^2}\right)\right).
              \]
              块级数收敛,而块内任意未完成部分的绝对值不超过 $(2k+1)/k^2\to0$,故原级数收敛.
        \item 题面中的 $(-1)^{\sqrt[3]n}$ 对非立方数不是实数,无法作为实级数讨论.若原题应为
              $(-1)^{[\sqrt[3]n]}$,则按 $k^3\leq n<(k+1)^3$ 分块,块和为
              \[
                  (-1)^k\left(\dfrac{3}{k}+O\left(\dfrac{1}{k^2}\right)\right),
              \]
              且块内余量趋于零,故同样收敛.指数中的取整符号需要核对原始资料.
    \end{enumerate}
\end{solution}

%例题8.3.3
\begin{example}
    讨论级数
    \[
        \sum\limits_{n=1}^{\infty}a_n
        =\dfrac{1}{1^p}-\dfrac{1}{2^q}+\dfrac{1}{3^p}-\dfrac{1}{4^q}+
        \cdots+\dfrac{1}{(2n-1)^p}-\dfrac{1}{(2n)^q}+\cdots,
    \]
    其中 $p>0,q>0$,的绝对收敛与条件收敛性.
\end{example}
\exampleblank

\begin{solution}
    绝对收敛等价于两个正项级数同时收敛,故恰在 $p>1$ 且 $q>1$ 时成立.

    若 $p=q=r$,两两分组后
    \[
        \dfrac{1}{(2n-1)^r}-\dfrac{1}{(2n)^r}=O(n^{-r-1}),
    \]
    所以原级数对一切 $r>0$ 均收敛;其中 $0<r\leq1$ 时为条件收敛.若
    $p\neq q$ 且至少一个不超过 $1$,分组通项最终由较小指数对应的项支配,部分和趋于 $+\infty$ 或 $-\infty$.因此条件收敛恰在 $p=q\leq1$ 时发生.
\end{solution}

%例题8.3.4
\begin{example}
    研究
    \[
        \sum\limits_{n=1}^{\infty}\dfrac{(-1)^{n-1}}{n^{p+\frac{1}{n}}}
    \]
    的绝对收敛与条件收敛性.
\end{example}
\exampleblank

\begin{solution}
    有
    \[
        n^{-p-\frac{1}{n}}=e^{-\frac{\ln n}{n}}n^{-p}\sim n^{-p}.
    \]
    当 $p>0$ 时该正数列最终递减至零,故原交错级数收敛;其绝对值级数恰在 $p>1$ 时收敛.因此 $p>1$ 时绝对收敛,$0<p\leq1$ 时条件收敛;$p\leq0$ 时通项不趋于零.
\end{solution}

%例题8.3.5
\begin{example}
    设 $f(x)$ 在区间 $\lbrack -1,1\rbrack$ 上有二阶连续导数,且 $\displaystyle\lim\limits_{x\to0}\dfrac{f(x)}{x}=0$,证明:
    \[
        \sum\limits_{n=1}^{\infty}f\left(\dfrac{1}{n}\right)
    \]
    绝对收敛.
\end{example}
\exampleblank

\begin{solution}
    题设极限先给出 $f(0)=0$,再由可导性得到 $f'(0)=0$.Taylor 公式给出:对每个 $n$,存在 $\xi_n$ 介于 $0$ 与 $1/n$ 之间,使
    \[
        f\left(\dfrac{1}{n}\right)=\dfrac{f''(\xi_n)}{2n^2}.
    \]
    由于 $f''$ 在紧区间上有界,故 $|f(1/n)|\leq C/n^2$,从而题中级数绝对收敛.
\end{solution}

%例题8.3.6
\begin{example}
    讨论级数
    \[
        \sum\limits_{n=1}^{\infty}\left(1+\dfrac{1}{2}+\dfrac{1}{3}+\cdots+\dfrac{1}{n}\right)\dfrac{\sin nx}{n}
    \]
    的敛散性.
\end{example}
\exampleblank

\begin{solution}
    记 $H_n=1+1/2+\cdots+1/n$.数列 $H_n/n$ 从某项起单调递减且趋于零.若 $x\notin2\pi\mathbb{Z}$,则
    \[
        \sum\limits_{k=1}^{N}\sin(kx)
        =\dfrac{\sin(Nx/2)\sin((N+1)x/2)}{\sin(x/2)}
    \]
    一致有界,故由 Dirichlet 判别法原级数收敛;若 $x\in2\pi\mathbb{Z}$,各项均为零.因此对每个实数 $x$,该级数都收敛.
\end{solution}

%例题8.3.7
\begin{example}
    设 $\displaystyle\sum\limits_{n=1}^{\infty}a_n$ 收敛,$\displaystyle\sum\limits_{n=1}^{\infty}(a_n+b_n)$ 绝对收敛,试证:$\displaystyle\sum\limits_{n=1}^{\infty}a_nb_n$ 也收敛.
\end{example}
\exampleblank

\begin{solution}
    当前命题不成立.取
    \[
        a_n=\dfrac{(-1)^{n-1}}{\sqrt n},
        \qquad b_n=-a_n.
    \]
    则 $\displaystyle \sum a_n$ 收敛,且 $\displaystyle \sum(a_n+b_n)=0$ 绝对收敛;但
    \[
        \sum\limits_{n=1}^{\infty}a_nb_n=-\sum\limits_{n=1}^{\infty}\dfrac{1}{n}
    \]
    发散.因此题目还缺少条件,需要核对原始资料.
\end{solution}

%例题8.3.8
\begin{example}
    判别下列级数的敛散性:
    \begin{enumerate}
        \item $\displaystyle\sum\limits_{n=1}^{\infty}\dfrac{(-1)^n}{\sqrt n+(-1)^{n-1}}$;
        \item $\displaystyle\sum\limits_{n=1}^{\infty}\ln\left(1+\dfrac{(-1)^n}{\sqrt[3]{n^2}}\right)$.
    \end{enumerate}
\end{example}
\exampleblank

\begin{solution}
    \begin{enumerate}
        \item 展开得
              \[
                  \dfrac{(-1)^n}{\sqrt n+(-1)^{n-1}}
                  =\dfrac{(-1)^n}{\sqrt n}+\dfrac{1}{n}+O(n^{-3/2}).
              \]
              第一项构成交错收敛级数,第二项为调和级数,故原级数发散到 $+\infty$.
        \item 令 $u_n=(-1)^nn^{-2/3}$,则
              \[
                  \ln(1+u_n)=u_n-\dfrac{1}{2}u_n^2+O(|u_n|^3).
              \]
              第一项形成收敛的交错级数,后两项绝对可和,故原级数收敛;但
              $|\ln(1+u_n)|\sim n^{-2/3}$,所以不是绝对收敛.
    \end{enumerate}
\end{solution}

%例题8.3.9
\begin{example}
    判断下列级数的敛散性:
    \begin{enumerate}
        \item $\displaystyle\sum\limits_{n=1}^{\infty}(-1)^n\left(1-\cos\dfrac{\pi}{\sqrt n}\right)$;
        \item $\displaystyle\sum\limits_{n=1}^{\infty}\sin\left(\pi\sqrt{n^2+1}\right)$.
    \end{enumerate}
\end{example}
\exampleblank

\begin{solution}
    \begin{enumerate}
        \item
              \[
                  1-\cos\dfrac{\pi}{\sqrt n}\sim\dfrac{\pi^2}{2n},
              \]
              故原级数与交错调和级数同型,条件收敛.
        \item 因
              \[
                  \sqrt{n^2+1}=n+\dfrac{1}{2n}+O(n^{-3}),
              \]
              所以
              \[
                  \sin\left(\pi\sqrt{n^2+1}\right)
                  =(-1)^n\left(\dfrac{\pi}{2n}+O(n^{-3})\right).
              \]
              因此该级数也条件收敛.
    \end{enumerate}
\end{solution}

%例题8.3.10
\begin{example}
    判断下列级数的敛散性:
    \begin{enumerate}
        \item $\displaystyle\sum\limits_{n=2}^{\infty}\dfrac{(-1)^n}{\sqrt n}
                  \left[(-1)^n+\sqrt n\sin\dfrac{1}{\sqrt n}\right]$;
        \item $\displaystyle\sum\limits_{n=1}^{\infty}(-1)^{\frac{n(n-1)}{2}}\dfrac{1}{\sqrt n}$.
    \end{enumerate}
\end{example}
\exampleblank

\begin{solution}
    \begin{enumerate}
        \item 通项化为
              \[
                  \dfrac{1}{\sqrt n}+(-1)^n\sin\dfrac{1}{\sqrt n}
                  =\dfrac{1+(-1)^n}{\sqrt n}+O(n^{-3/2}).
              \]
              偶数项的主项形成发散的正项级数,而误差绝对可和,故原级数发散到 $+\infty$.
        \item 系数 $(-1)^{n(n-1)/2}$ 以 $+,-,-,+$ 为周期,其部分和有界;又 $n^{-1/2}$ 单调趋零,故由 Dirichlet 判别法收敛.绝对值级数发散,所以是条件收敛.
    \end{enumerate}
\end{solution}

记
\[
    a_n^+=\begin{cases}a_n, & a_n\geq0, \\0,&a_n<0,
    \end{cases}
    \qquad
    a_n^-=\begin{cases}-a_n, & a_n\leq0, \\0,&a_n>0,
    \end{cases}
\]
则 $a_n^+,a_n^-\geq0$ 且 $a_n=a_n^+-a_n^-$,$|a_n|=a_n^++a_n^-$.

\begin{theorem}[\markstar]
    设 $\displaystyle\sum\limits_{n=1}^{\infty}a_n$ 绝对收敛,则:
    \begin{enumerate}
        \item $\displaystyle\sum\limits_{n=1}^{\infty}a_n^+$,$\displaystyle\sum\limits_{n=1}^{\infty}a_n^-$ 收敛;
        \item $\displaystyle\sum\limits_{n=1}^{\infty}a_n$ 收敛.
    \end{enumerate}
\end{theorem}

\begin{corollary}[\markstar]
    若 $\displaystyle\sum\limits_{n=1}^{\infty}a_n$ 条件收敛,则
    \[
        \sum\limits_{n=1}^{\infty}a_n^+=+\infty,
        \qquad
        \sum\limits_{n=1}^{\infty}a_n^-=+\infty.
    \]
\end{corollary}

\begin{remark}
    由此可得条件收敛级数 $\displaystyle\sum\limits_{n=1}^{\infty}a_n$ 收敛的原因在于其中正项与负项大部分抵消了.
\end{remark}

%例题8.3.11
\begin{example}
    判别级数 $\displaystyle I=\sum\limits_{n=1}^{\infty}a_n$ 的收敛性:
    \begin{enumerate}
        \item $\displaystyle a_n=\dfrac{(-1)^n}{\sqrt[3]{n}+(-1)^{n-1}}$;
        \item $\displaystyle a_n=\ln\left(1+\dfrac{(-1)^{n-1}}{(n+1)^\alpha}\right)$;
        \item $\displaystyle a_n=\dfrac{(-1)^n}{[2n+(-1)^n]^\alpha}$;
        \item $\displaystyle a_n=\dfrac{(-1)^{n-1}}{(\sqrt{n+1}+(-1)^n)^\alpha}$.
    \end{enumerate}
\end{example}
\exampleblank

\begin{solution}
    逐项展开:
    \begin{enumerate}
        \item
              \[
                  a_n=\dfrac{(-1)^n}{n^{1/3}}+\dfrac{1}{n^{2/3}}+O(n^{-1}),
              \]
              故发散.
        \item 当 $\alpha>0$ 时,
              \[
                  a_n=\dfrac{(-1)^{n-1}}{(n+1)^\alpha}
                  -\dfrac{1}{2(n+1)^{2\alpha}}+O(n^{-3\alpha}).
              \]
              因而 $\alpha>1/2$ 时收敛,$0<\alpha\leq1/2$ 时发散;其中
              $\alpha>1$ 绝对收敛,$1/2<\alpha\leq1$ 条件收敛.
        \item
              \[
                  a_n=\dfrac{(-1)^n}{(2n)^\alpha}
                  -\dfrac{\alpha}{(2n)^{\alpha+1}}+O(n^{-\alpha-2}).
              \]
              故 $\alpha>0$ 时收敛,且 $\alpha>1$ 时绝对收敛,$0<\alpha\leq1$ 时条件收敛;$\alpha\leq0$ 时通项不趋于零.
        \item
              \[
                  a_n=(-1)^{n-1}n^{-\alpha/2}
                  +\alpha n^{-(\alpha+1)/2}+O(n^{-(\alpha+2)/2}).
              \]
              故恰在 $\alpha>1$ 时收敛;$\alpha>2$ 时绝对收敛,$1<\alpha\leq2$ 时条件收敛.
    \end{enumerate}
\end{solution}

%例题8.3.12
\begin{example}
    \begin{enumerate}
        \item 设 $x\neq k\pi$($k\in\mathbb{Z}$),讨论 $\displaystyle\sum\limits_{n=1}^{\infty}\dfrac{\cos nx}{n^p}$ 与 $\displaystyle\sum\limits_{n=1}^{\infty}\dfrac{\sin nx}{n^p}$ 的条件收敛性与绝对收敛性;
        \item 讨论级数 $\displaystyle\sum\limits_{n=1}^{\infty}\dfrac{(-1)^n\sin n}{n}$ 与 $\displaystyle\sum\limits_{n=1}^{\infty}\dfrac{(-1)^n\sin^2n}{n}$ 的条件收敛性与绝对收敛性.
    \end{enumerate}
\end{example}
\exampleblank

\begin{solution}
    \begin{enumerate}
        \item 当 $p>0$ 时,$n^{-p}$ 单调趋零,而 $\displaystyle \sum\cos(nx)$,$\displaystyle \sum\sin(nx)$ 的部分和在 $x\notin\pi\mathbb{Z}$ 时有界,故两级数收敛.它们在 $p>1$ 时绝对收敛,在 $0<p\leq1$ 时条件收敛;$p\leq0$ 时一般不收敛.
        \item 因
              \[
                  (-1)^n\sin n=\sin(n(1+\pi)),
              \]
              第一个级数由 Dirichlet 判别收敛;其绝对值级数发散.又
              \[
                  (-1)^n\sin^2n=\dfrac{(-1)^n}{2}-\dfrac{(-1)^n\cos2n}{2},
              \]
              第二个级数也由 Dirichlet 判别收敛,绝对值级数仍发散.因此二者均条件收敛.
    \end{enumerate}
\end{solution}

关于级数重排的问题,有以下两个结论.

\begin{theorem}[\markstar]
    交换绝对收敛级数中无穷多项的次序,所得的新级数仍然绝对收敛,其和也不变.
\end{theorem}

\begin{theorem}[\markstar\ Riemann 定理]
    若 $\displaystyle\sum\limits_{n=1}^{\infty}a_n$ 条件收敛,则适当交换各项的次序,可使其收敛到任意事先指定的实数 $S$,也可以使其发散到 $+\infty$ 或 $-\infty$.
\end{theorem}


\newpage

\section{级数的应用 *}


\subsection{与极限的联系}

%例题8.4.1
\begin{example}
    设
    \[
        x_n=\dfrac{x^n}{n!\,3^n},
    \]
    求 $\displaystyle\lim\limits_{n\to\infty}x_n$.
\end{example}
\exampleblank

\begin{solution}
    固定 $x$.由
    \[
        \left|\dfrac{x_{n+1}}{x_n}\right|=\dfrac{|x|}{3(n+1)}\longrightarrow0,
    \]
    可知 $x_n\to\boxed0$.
\end{solution}

%例题8.4.2
\begin{example}
    设 $a_n>0$,$\displaystyle\sum\limits_{n=1}^{\infty}a_n=1$,$\displaystyle A_n=\sum\limits_{k=1}^na_k$,求极限
    \[
        \lim\limits_{n\to\infty}\dfrac{e^{A_n}-e^{A_{n-1}}}{A_n^e-A_{n-1}^e}.
    \]
\end{example}
\exampleblank

\begin{solution}
    由 Cauchy 中值定理,存在介于 $A_{n-1}$ 与 $A_n$ 之间的 $\xi_n,\eta_n$,使
    \[
        \dfrac{e^{A_n}-e^{A_{n-1}}}{A_n^e-A_{n-1}^e}
        =\dfrac{e^{\xi_n}}{e\eta_n^{e-1}}.
    \]
    因 $A_n\to1$,有 $\xi_n,\eta_n\to1$,故所求极限为 $\boxed1$.
\end{solution}

%例题8.4.3
\begin{example}
    设数列 $\{a_n\}$ 满足条件:
    \begin{enumerate}
        \item $0<a_k\leq100a_n$($n=k+1,k+2,\ldots$);
        \item $\displaystyle\sum\limits_{n=1}^{\infty}a_n$ 收敛.
    \end{enumerate}
    求证:$\displaystyle\lim\limits_{n\to\infty}na_n=0$.
\end{example}
\exampleblank

\begin{solution}
    固定 $n$,对 $m=n+1,\ldots,2n$ 使用条件 (1),有 $a_n\leq100a_m$.故
    \[
        na_n\leq100\sum\limits_{m=n+1}^{2n}a_m.
    \]
    右端是收敛级数的尾和,趋于零,所以 $na_n\to0$.
\end{solution}

%例题8.4.4
\begin{example}[\markstar]
    设正项级数 $\displaystyle\sum\limits_{n=1}^{\infty}a_n$ 收敛,试证:
    \[
        \lim\limits_{n\to\infty}\dfrac{\sum\limits_{k=1}^nka_k}{n}=0.
    \]
\end{example}
\exampleblank

\begin{solution}
    任取 $N<n$,分拆为
    \[
        \dfrac{1}{n}\sum\limits_{k=1}^{n}ka_k
        =\dfrac{1}{n}\sum\limits_{k=1}^{N}ka_k
        +\dfrac{1}{n}\sum\limits_{k=N+1}^{n}ka_k.
    \]
    第一项在 $n\to\infty$ 时趋于零,第二项不超过
    $\displaystyle \sum_{k=N+1}^{\infty}a_k$.再令 $N\to\infty$,即得极限为零.
\end{solution}

%例题8.4.5
\begin{example}
    设 $\displaystyle S_n=\sum\limits_{k=1}^n\dfrac{\sin k}{k}=S_n^++S_n^-$,其中 $S_n^+$ 和 $S_n^-$ 分别是正项之和和负项之和.证明:$\displaystyle\lim\limits_{n\to\infty}\dfrac{S_n^+}{S_n^-}$ 存在并求其值.
\end{example}
\exampleblank

\begin{solution}
    由 Dirichlet 判别,$S_n$ 收敛;而 $\displaystyle \sum|\sin n|/n$ 发散,故其正项和
    $S_n^+\to+\infty$,负项和 $S_n^-\to-\infty$.由于
    $S_n^++S_n^-=S_n=O(1)$,有
    \[
        \dfrac{S_n^+}{S_n^-}
        =-1+\dfrac{S_n}{S_n^-}\longrightarrow\boxed{-1}.
    \]
\end{solution}



\subsection{Cesaro 和}

\begin{definition}
    设 $\displaystyle\sum\limits_{n=1}^{\infty}a_n$ 为一无穷级数,记
    \[
        \sigma_n=\dfrac{S_1+S_2+\cdots+S_n}{n},
    \]
    其中 $S_n$ 为前 $n$ 项和.若 $\{\sigma_n\}$ 收敛,即 $\displaystyle\lim\limits_{n\to\infty}\sigma_n=\sigma$,则称 $\displaystyle\sum\limits_{n=1}^{\infty}a_n$ 在 Cesaro 意义下收敛,$\sigma$ 称为该级数的 Cesaro 和,记为
    \[
        \sum\limits_{n=1}^{\infty}a_n=\sigma(C).
    \]
    此时称级数可以 Cesaro 求和.
\end{definition}

\begin{remark}
    若 $\displaystyle\sum\limits_{n=1}^{\infty}a_n$ 收敛,即 $\displaystyle\lim\limits_{n\to\infty}S_n=S$,则 $\displaystyle\lim\limits_{n\to\infty}\sigma_n=S$.由此可以看出 Cesaro 和为更广义的和,能保证原来收敛的级数仍可求 Cesaro 和,部分发散的级数在 Cesaro 意义下亦可求和,例如
    \[
        \sum\limits_{n=1}^{\infty}(-1)^{n-1}=\dfrac{1}{2}(C).
    \]
\end{remark}

%例题8.4.6
\begin{example}
    求下列级数的 Cesaro 和:
    \begin{enumerate}
        \item $1+0-1+1+0-1+\cdots$;
        \item $\sin x+\sin2x+\sin3x+\cdots$,其中 $0<x<2\pi$.
    \end{enumerate}
\end{example}
\exampleblank

\begin{solution}
    \begin{enumerate}
        \item 部分和以 $1,1,0$ 为周期,故其算术平均趋于 $\boxed{2/3}$.
        \item 由
              \[
                  \sum\limits_{k=1}^{N}\sin(kx)
                  =\dfrac{\cos(x/2)-\cos((N+1/2)x)}{2\sin(x/2)},
              \]
              对 $N=1,ldots,n$ 再取平均.振荡余项的平均趋于零,故 Cesaro 和为
              \[
                  \boxed{\dfrac{1}{2}\cot\dfrac{x}{2}}.
              \]
    \end{enumerate}
\end{solution}

%例题8.4.7
\begin{example}[\markstar\ 可以 Cesaro 求和的必要条件]
    设 $\displaystyle\sum\limits_{n=1}^{\infty}a_n$ 可以 Cesaro 求和,则
    \[
        \lim\limits_{n\to\infty}\dfrac{S_n}{n}=0,
        \qquad
        \lim\limits_{n\to\infty}\dfrac{a_n}{n}=0.
    \]
\end{example}
\exampleblank

\begin{solution}
    记 $\sigma_n=(S_1+\cdots+S_n)/n\to s$.则
    \[
        \dfrac{S_n}{n}=\sigma_n-\dfrac{n-1}{n}\sigma_{n-1}\longrightarrow0.
    \]
    又 $a_n=S_n-S_{n-1}$,所以
    \[
        \dfrac{a_n}{n}=\dfrac{S_n}{n}-\dfrac{n-1}{n}\dfrac{S_{n-1}}{n-1}\longrightarrow0.
    \]
\end{solution}

%例题8.4.8
\begin{example}[\markstar]
    设 $\displaystyle\lim\limits_{n\to\infty}na_n=0$,则
    \[
        \sum\limits_{n=1}^{\infty}a_n\text{ 在通常意义下收敛}
        \quad\Longleftrightarrow\quad
        \sum\limits_{n=1}^{\infty}a_n\text{ 可以 Cesaro 求和}.
    \]
\end{example}
\exampleblank

\begin{solution}
    通常收敛必然蕴含 Cesaro 可和.反之设 Cesaro 和为 $s$.恒等式
    \[
        S_n-\sigma_n=\dfrac{1}{n}\sum\limits_{k=1}^{n}(k-1)a_k
    \]
    成立.由 $ka_k\to0$ 及 Cesaro 平均定理,右端趋于零,故
    $S_n\to s$,即通常收敛.
\end{solution}

%例题8.4.9
\begin{example}[\markstar]
    设 $\displaystyle\sum\limits_{n=1}^{\infty}a_n$ 可以 Cesaro 求和,证明:
    \[
        \sum\limits_{n=1}^{\infty}a_n\text{ 在通常意义下收敛}
        \quad\Longleftrightarrow\quad
        \lim\limits_{n\to\infty}\dfrac{a_1+2a_2+\cdots+na_n}{n}=0.
    \]
\end{example}
\exampleblank

\begin{solution}
    仍记 Cesaro 均值为 $\sigma_n\to s$.由
    \[
        S_n-\sigma_n=\dfrac{1}{n}\sum\limits_{k=1}^{n}(k-1)a_k
    \]
    可知,若题中加权平均趋于零,则 $S_n\to s$.反之,若 $S_n$ 通常收敛,则由 Kronecker 引理
    \[
        \dfrac{1}{n}\sum\limits_{k=1}^{n}ka_k\longrightarrow0.
    \]
    这就证明了等价性.
\end{solution}



\subsection{Cauchy 乘积}

\begin{definition}
    级数 $\displaystyle\sum\limits_{n=1}^{\infty}a_n$ 与 $\displaystyle\sum\limits_{n=1}^{\infty}b_n$ 的乘积定义为 $\displaystyle\sum\limits_{n=1}^{\infty}c_n$,其中
    \[
        c_n=a_1b_n+a_2b_{n-1}+\cdots+a_nb_1.
    \]
    这种级数的乘积称为 Cauchy 乘积,又称为对角线乘积.
\end{definition}

\begin{theorem}[\markstar\ Cauchy]
    设 $\displaystyle\sum\limits_{n=1}^{\infty}a_n$ 与 $\displaystyle\sum\limits_{n=1}^{\infty}b_n$ 都绝对收敛,它们的和分别为 $A,B$,则将所有 $a_ib_j$ 按任何方式求和而成的级数也绝对收敛,且和为 $AB$.
\end{theorem}

\begin{theorem}[\markstar\ Mertens]
    设 $\displaystyle\sum\limits_{n=1}^{\infty}a_n$ 与 $\displaystyle\sum\limits_{n=1}^{\infty}b_n$ 收敛且至少有一个绝对收敛,它们的和分别为 $A,B$,则它们的 Cauchy 乘积 $\displaystyle\sum\limits_{n=1}^{\infty}c_n$ 收敛,且和为 $AB$.
\end{theorem}

\begin{theorem}[\markstar\ Abel]
    设 $\displaystyle\sum\limits_{n=1}^{\infty}a_n$,$\displaystyle\sum\limits_{n=1}^{\infty}b_n$ 与它们的 Cauchy 乘积 $\displaystyle\sum\limits_{n=1}^{\infty}c_n$ 都收敛,且和分别为 $A,B,C$,则 $C=AB$.
\end{theorem}

%例题8.4.10
\begin{example}[\markstar]
    设 $|x|<1$,则
    \[
        \dfrac{1}{1-x}\sum\limits_{n=0}^{\infty}a_nx^n
        =\left(\sum\limits_{n=0}^{\infty}x^n\right)
        \left(\sum\limits_{n=0}^{\infty}a_nx^n\right)
        =\sum\limits_{n=0}^{\infty}(a_0+a_1+\cdots+a_n)x^n.
    \]
    特别地,
    \[
        \left(\dfrac{1}{1-x}\right)^2
        =\dfrac{1}{1-x}\sum\limits_{n=0}^{\infty}x^n
        =\sum\limits_{n=0}^{\infty}(n+1)x^n.
    \]
\end{example}
\exampleblank

\begin{solution}
    两级数在 $|x|<1$ 时绝对收敛,故可作 Cauchy 乘积.$x^n$ 的系数为
    \[
        \sum\limits_{k=0}^{n}a_k=a_0+a_1+\cdots+a_n,
    \]
    从而得到第一式.再取 $a_n=1$,系数即为 $n+1$,得到第二式.
\end{solution}

%例题8.4.11
\begin{example}
    证明:$\displaystyle\sum\limits_{n=1}^{\infty}\dfrac{a^n}{n!}$ 与 $\displaystyle\sum\limits_{n=1}^{\infty}\dfrac{b^n}{n!}$ 绝对收敛,且它们的 Cauchy 乘积等于
    \[
        \sum\limits_{n=0}^{\infty}\dfrac{(a+b)^n}{n!}.
    \]
\end{example}
\exampleblank

\begin{solution}
    比式判别给出两个指数级数均绝对收敛.其 Cauchy 乘积中 $n$ 次项系数为
    \[
        \sum\limits_{k=0}^{n}\dfrac{a^k}{k!}\dfrac{b^{n-k}}{(n-k)!}
        =\dfrac{1}{n!}\sum\limits_{k=0}^{n}C_n^ka^kb^{n-k}
        =\dfrac{(a+b)^n}{n!}.
    \]
    故乘积正是题中级数.
\end{solution}

%例题8.4.12
\begin{example}
    求 $\displaystyle\sum\limits_{n=1}^{\infty}nx^{n-1}$ 与 $\displaystyle\sum\limits_{n=1}^{\infty}(-1)^{n-1}nx^{n-1}$ 的 Cauchy 乘积.
\end{example}
\exampleblank

\begin{solution}
    当 $|x|<1$ 时,两级数分别为
    \[
        \dfrac{1}{(1-x)^2},
        \qquad
        \dfrac{1}{(1+x)^2}.
    \]
    故 Cauchy 乘积为
    \[
        \dfrac{1}{(1-x^2)^2}
        =\sum\limits_{m=0}^{\infty}(m+1)x^{2m}.
    \]
    即奇次幂系数为零,$x^{2m}$ 的系数为 $m+1$.
\end{solution}

%例题8.4.13
\begin{example}
    求下列级数 $S=\sum c_n$ 的和:
    \begin{enumerate}
        \item $\displaystyle\sum\limits_{n=0}^{\infty}c_n$,其中 $\displaystyle c_n=\sum\limits_{k=0}^nx^ky^{n-k}$,$|x|<1,|y|<1$;
        \item $\displaystyle\sum\limits_{n=1}^{\infty}c_n$,其中 $\displaystyle c_n=\sum\limits_{k=1}^{\infty}\dfrac{1}{k(k+1)(n-k+1)!}$.
    \end{enumerate}
\end{example}
\exampleblank

\begin{solution}
    \begin{enumerate}
        \item 这是两个几何级数的 Cauchy 乘积,故
              \[
                  S=\left(\sum\limits_{k=0}^{\infty}x^k\right)
                  \left(\sum\limits_{j=0}^{\infty}y^j\right)
                  =\boxed{\dfrac{1}{(1-x)(1-y)}}.
              \]
        \item 当前题面令 $k$ 从 $1$ 到 $\infty$,但当 $k>n+1$ 时阶乘
              $(n-k+1)!$ 无定义.若自然的上限应为 $k=n$,则令 $j=n-k+1$,交换绝对收敛求和次序可得
              \[
                  S=\left(\sum\limits_{k=1}^{\infty}\dfrac{1}{k(k+1)}\right)
                  \left(\sum\limits_{j=1}^{\infty}\dfrac{1}{j!}\right)
                  =\boxed{e-1}.
              \]
              原题的求和上限需要核对.
    \end{enumerate}
\end{solution}



\subsection{无穷乘积}

\begin{definition}
    设 $\displaystyle\prod\limits_{n=1}^{\infty}p_n$ 是一个无穷乘积,称
    \[
        P_n=\prod\limits_{k=1}^np_k=p_1\cdots p_n
    \]
    为这个无穷乘积的部分乘积.若 $\displaystyle\lim\limits_{n\to\infty}P_n=p$(有限且 $\neq0$),则称这个无穷乘积是收敛的,反之称它是发散的.
\end{definition}

为什么要去掉 $p=0$ 的情况呢?实则是为了得到类似级数收敛的必要条件.

\begin{theorem}[\markstar]
    无穷乘积 $\displaystyle\prod\limits_{n=1}^{\infty}p_n$ 收敛的必要条件是
    \[
        \lim\limits_{n\to\infty}p_n=1.
    \]
\end{theorem}

\begin{theorem}[\markstar]
    \[
        \prod\limits_{n=1}^{\infty}(1+a_n)\text{ 收敛}
        \quad\Longleftrightarrow\quad
        \sum\limits_{n=1}^{\infty}\ln(1+a_n)\text{ 收敛}.
    \]
\end{theorem}

\begin{remark}
    收敛情况下,若后者和为 $S$,则前者积为 $e^S$.
\end{remark}

\begin{theorem}[\markstar]
    设 $a_n>0$($n>n_0$),则 $\displaystyle\prod\limits_{n=1}^{\infty}(1+a_n)$ 与 $\displaystyle\sum\limits_{n=1}^{\infty}a_n$ 同敛散.
\end{theorem}

\begin{theorem}[\markstar]
    如果 $\displaystyle\sum\limits_{n=1}^{\infty}a_n^2<+\infty$,那么 $\displaystyle\prod\limits_{n=1}^{\infty}(1+a_n)$ 与 $\displaystyle\sum\limits_{n=1}^{\infty}a_n$ 同敛散.
\end{theorem}

\begin{theorem}
    设 $-1<a_n<0$,且 $\displaystyle\sum\limits_{n=1}^{\infty}a_n$ 发散,则 $\displaystyle\prod\limits_{n=1}^{\infty}(1+a_n)$ 发散到 $0$.
\end{theorem}

\begin{theorem}
    如果 $\displaystyle\sum\limits_{n=1}^{\infty}a_n$ 收敛,设 $\displaystyle\sum\limits_{n=1}^{\infty}a_n^2$ 发散,那么 $\displaystyle\prod\limits_{n=1}^{\infty}(1+a_n)$ 发散到 $0$.
\end{theorem}

%例题8.4.14
\begin{example}
    证明下列等式:
    \begin{enumerate}
        \item $\displaystyle\prod\limits_{n=2}^{\infty}\dfrac{n^3-1}{n^3+1}=\dfrac{2}{3}$;
        \item $\displaystyle\prod\limits_{n=0}^{\infty}\left[1+\left(\dfrac{1}{2}\right)^{2^n}\right]=2$.
    \end{enumerate}
\end{example}
\exampleblank

\begin{solution}
    \begin{enumerate}
        \item 前 $N$ 项乘积为
              \[
                  \begin{aligned}
                      \prod\limits_{n=2}^{N}\dfrac{n^3-1}{n^3+1}
                       & =\prod\limits_{n=2}^{N}\dfrac{n-1}{n+1}
                      \prod\limits_{n=2}^{N}\dfrac{n^2+n+1}{n^2-n+1} \\
                       & =\dfrac{2}{N(N+1)}\dfrac{N^2+N+1}{3}
                      \longrightarrow\dfrac{2}{3}.
                  \end{aligned}
              \]
        \item 恒等式
              \[
                  \prod\limits_{n=0}^{N}(1+x^{2^n})=\dfrac{1-x^{2^{N+1}}}{1-x}
              \]
              在 $x=1/2$ 时给出极限 $2$.
    \end{enumerate}
\end{solution}

%例题8.4.15
\begin{example}
    讨论下列无穷乘积的敛散性:
    \begin{enumerate}
        \item $\displaystyle\prod\limits_{n=1}^{\infty}\dfrac{1}{n}$;
        \item $\displaystyle\prod\limits_{n=1}^{\infty}\dfrac{(n+1)^2}{n(n+2)}$;
        \item $\displaystyle\prod\limits_{n=1}^{\infty}\sqrt[n]{1+\dfrac{1}{n}}$.
    \end{enumerate}
\end{example}
\exampleblank

\begin{solution}
    \begin{enumerate}
        \item 部分乘积为 $1/N!\to0$,故按非零有限极限的定义,该无穷乘积发散.
        \item 裂项相消得
              \[
                  \prod\limits_{n=1}^{N}\dfrac{(n+1)^2}{n(n+2)}
                  =\dfrac{2(N+1)}{N+2}\longrightarrow2,
              \]
              故收敛.
        \item 取对数后,
              \[
                  \sum\limits_{n=1}^{\infty}\dfrac{1}{n}\ln\left(1+\dfrac{1}{n}\right)
              \]
              的通项与 $1/n^2$ 等价,故对数级数收敛,原无穷乘积收敛到正数.
    \end{enumerate}
\end{solution}

%例题8.4.16
\begin{example}
    讨论下列无穷乘积的敛散性:
    \begin{enumerate}
        \item $\displaystyle\prod\limits_{n=1}^{\infty}\dfrac{n}{\sqrt{n^2+1}}$;
        \item $\displaystyle\prod\limits_{n=2}^{\infty}\left(\dfrac{n^2-1}{n^2+1}\right)^p\quad(p\in\mathbb{R})$.
    \end{enumerate}
\end{example}
\exampleblank

\begin{solution}
    \begin{enumerate}
        \item 取对数后通项为
              \[
                  -\dfrac{1}{2}\ln\left(1+\dfrac{1}{n^2}\right)\sim-\dfrac{1}{2n^2},
              \]
              故对数级数收敛,原乘积收敛到正数.
        \item 对任意固定 $p\in\mathbb{R}$,
              \[
                  p\ln\dfrac{n^2-1}{n^2+1}
                  =-\dfrac{2p}{n^2}+O(n^{-4}).
              \]
              对数级数绝对收敛,故原无穷乘积对每个实数 $p$ 都收敛到正数($p=0$ 时等于 $1$).
    \end{enumerate}
\end{solution}
